\documentclass[10pt,conference]{IEEEtran}
\IEEEoverridecommandlockouts
\usepackage[utf8]{inputenc}
\usepackage[T1]{fontenc}
\usepackage{microtype}
\usepackage{amsmath,amssymb}
\usepackage{amsthm}
\usepackage{graphicx}
\usepackage{hyperref}
\usepackage{booktabs}
\usepackage{multirow}
\usepackage{url}
\usepackage{algorithm}
\usepackage[table]{xcolor}
\usepackage{algpseudocode}
\usepackage{threeparttable}
\usepackage{adjustbox}
\usepackage{cite}
\usepackage{setspace}
\usepackage{subcaption}
\usepackage{listings}
\usepackage{helvet}
\usepackage{tikz}
\usetikzlibrary{positioning, arrows.meta, shapes.geometric, calc, fit, backgrounds}

\newtheorem{definition}{Definition}

% Listings style
\lstset{
    basicstyle=\fontfamily{phv}\fontsize{5.5pt}{6pt}\selectfont,
    frame=single,
    backgroundcolor=\color{yellow!10},
    showstringspaces=false,
    tabsize=1,
    breaklines=true,
    captionpos=b,
    framesep=2pt,
    emph={
        weiRaised,
        token.allocate,
        totalRaised,
        msg.sender.transfer,
        multisig.transfer
    },
    emphstyle=\color{black}\bfseries
}
\onehalfspacing

% Bibliography style
\bibliographystyle{IEEEtran}

% Title and Authors
\title{DeepEdu-v1: Efficient and Scalable Agentic LLMs for Vietnamese Education}
\author{
    % \IEEEauthorblockN{Quang Nguyen, Hieu Nguyen, Hien Hoang, Toan Pham, Cong Tran\IEEEauthorrefmark{1}, Nam Vu\IEEEauthorrefmark{1}}
    \IEEEauthorblockN{Quang Nguyen, Hieu Nguyen, Hien Hoang, Toan Pham, Cong Tran, Nam Vu}
    \IEEEauthorblockA{
        \textit{Posts and Telecommunications Institute of Technology, Ha Noi, Vietnam} \\
        Email: \{QuangNM.B21CN629, HieuNPT.B23CE031, HienHC.B23CE028, \\ ToanPQ.B23CN835\}@stu.ptit.edu.vn, congtt@ptit.edu.vn
    }
    % \thanks{\IEEEauthorrefmark{1}Corresponding author: Cong Tran, Nam Vu}
}

\begin{document}

\maketitle

% Activates the IEEEtranBSTCTL control entry (truncates long author lists to et al.)
\bstctlcite{IEEEexample:BSTcontrol}

\begin{abstract}
AI tutoring could markedly improve learning outcomes for students in developing regions such as Vietnam, yet the two obvious paths both fall short. Cloud assistants such as ChatGPT route sensitive student data to foreign servers---violating data-sovereignty laws such as Vietnam's Decree 53---and, pre-trained on Western-centric corpora, are not organized around the national textbook curriculum, so their knowledge of local content is unsystematic and frequently hallucinated. Self-hosting an open model keeps data on-premise but hits a two-fold wall: post-training quantization (AWQ, GPTQ) tames the static weight footprint, yet the dynamic KV cache and prefill latency of long tutoring contexts still cause out-of-memory failures and slow responses on consumer GPUs, while the model keeps hallucinating on region-specific material.

We present \textbf{DeepEdu-v1}, an AI-tutoring system for Vietnamese education built on \textbf{SCALE} (Self-improving Context-Aware Learning Engine), a framework with two innovations. First, a long-context inference engine amortizes token selection from per-sub-chunk to per-cluster granularity; on long-context retrieval it issues $7.7\times$ fewer retrieval calls than a state-of-the-art selective-attention baseline, cutting prefill latency (TTFT) by roughly $35\%$ while matching or improving task accuracy. Second, a self-improving agentic layer continuously curates a verified playbook from past interactions instead of fine-tuning, a design intended to progressively reduce reliance on dominant-language priors as trustworthy local knowledge accumulates. In its deployed configuration, DeepEdu achieves a nearly $2\times$ TTFT speedup over standard vLLM serving and lifts agentic accuracy from $70.0\%$ to $79.5\%$ on complex tasks, with the strongest per-track gains across financial-reasoning and interactive-agent benchmarks. Together, these results indicate that scalable, self-improving, curriculum-grounded AI tutoring can be democratized in resource-constrained regions.
\end{abstract}

\begin{IEEEkeywords}
Large Language Models, AI Education, Long-Context Inference, Model Quantization, Data Sovereignty, Self-improving Agents
\end{IEEEkeywords}

\section{Introduction}

Large Language Models (LLMs) have moved beyond text generation into multimodal, tool-augmented agents capable of complex reasoning, planning, and personalized interaction \cite{touvron2023llama, bai2023qwen, yao2023react}. In education, this puts within reach something long unavailable in resource-constrained systems: a one-on-one tutor that adapts to each student's pace and cognitive profile \cite{llmeducation2026}. Our goal is \textbf{DeepEdu}: an AI-tutoring platform that gives Vietnamese students a more effective, curriculum-aligned learning path, with LLM agents as its foundation. Realizing this vision, however, runs into two obstacles that a generic chatbot does not solve.

A natural first question is why not simply use a cloud assistant such as ChatGPT. Two reasons rule it out. First, \emph{data sovereignty}: educational records contain highly sensitive personally identifiable information, and routing them to foreign servers directly violates national data-localization laws such as Vietnam's Decree 53 \cite{vietnam_decree53}, compelling institutions to self-host open models strictly on-premise. Second, \emph{curriculum grounding}: pre-trained on English-dominant, Western-centric corpora \cite{cahyawijaya2024llms}, such models are not organized around the Vietnamese textbook curriculum (\emph{Sach Giao Khoa}, SGK); their knowledge of local content is unsystematic and frequently hallucinated \cite{vnexpress_chatgpt_education_2023}.\footnote{For instance, early deployments of ChatGPT in Vietnam generated severe historical hallucinations, such as claiming national leaders were founders of the World Health Organization (WHO), or confidently asserting that Emperor Quang Trung and his birth name Nguyen Hue were two distinct kings from opposing dynasties \cite{dnto_chatgpt_2023}. Such algorithmic hallucinations highlight the severe pedagogical risks of deploying ungrounded models in regional education.} A trustworthy tutor must therefore be both self-hosted and explicitly grounded in localized curricular knowledge.

Both requirements are hard on the infrastructure available to public schools. Self-hosting on limited hardware creates a two-fold bottleneck. The static footprint of model weights can be tamed by post-training quantization such as AWQ \cite{lin2024awq} or GPTQ \cite{frantar2023gptq}; but as a tutoring session stretches toward a 1M-token context, the \emph{dynamic} KV cache surpasses the static footprint and, together with quadratic prefill latency, causes Out-Of-Memory (OOM) failures and sluggish responses on consumer-grade GPUs---a bottleneck that I/O-aware kernels such as FlashAttention \cite{dao2022flashattention} and PagedAttention \cite{kwon2023efficient} alleviate but do not remove. Grounding is equally costly the obvious way: continually fine-tuning on evolving curricula is computationally prohibitive and risks catastrophic forgetting \cite{efficientcpt2025naacl}.

We address both obstacles with \textbf{DeepEdu-v1}, built on the \textbf{SCALE} framework (\textbf{S}elf-improving \textbf{C}ontext-\textbf{A}ware \textbf{L}earning \textbf{E}ngine). SCALE pairs an efficient long-context inference engine---which amortizes selective sparse attention across similar query chunks to keep latency low as context grows---with a self-improving agentic layer that curates a verified, curriculum-grounded playbook from past interactions rather than updating weights. The main contributions of our work are as follows:

\begin{enumerate}
    \item We introduce \textit{Similarity Chunk Rolling}, a long-context inference engine that amortizes selective sparse attention from per-sub-chunk to per-cluster granularity. By issuing a single retrieval per cluster, it reduces retrieval invocations by $7.7\times$ and prefill latency by roughly $35\%$ over a state-of-the-art selective-attention baseline (TokenSelect), while preserving---and on structured-retrieval tasks improving---task accuracy.
    \item We design a self-improving agentic layer for DeepEdu that, through query-aware context engineering, curates a verified playbook grounding the model in local knowledge and progressively reducing reliance on dominant-language priors---without computationally expensive weight updates.
    \item We show, across grounded financial-reasoning and interactive-agent benchmarks, that this self-improving loop delivers the strongest per-track gains over a strong agentic baseline, and that in its deployed configuration DeepEdu sustains a nearly $2\times$ TTFT speedup---a step toward democratizing localizable, self-improving AI tutoring in resource-constrained regions under data-sovereignty constraints.
\end{enumerate}

\section{Related Work}
\label{sec:related}
To contextualize the contributions of DeepEdu, this section surveys
recent advancements across the architectural evolution of LLMs and 
their application to education, the principal techniques for 
compressing models and optimizing long-context inference under hardware 
constraints, the challenges of deploying LLMs on low-resource languages 
and within localized cultural contexts, and the emerging paradigm of 
self-improving agentic frameworks.

\subsection{The Evolution of Large Language Models}
The landscape of Natural Language Processing (NLP) has been fundamentally transformed by the scaling of neural architectures, passing through several critical milestones. This trajectory is rooted in the Transformer architecture and early scaling laws \cite{vaswani2017attention, brown2020language}. Subsequently, the community witnessed the rise of powerful open-weights foundation models that democratized high-performance language modeling \cite{touvron2023llama, bai2023qwen}. To optimize the intrinsic trade-off between model capacity and inference cost, recent architectures have increasingly adopted the Mixture-of-Experts (MoE) paradigm \cite{yang2024qwen2, qwen252024, jiang2024mixtral}. This architectural evolution has recently culminated in highly efficient, post-training reinforced paradigms (such as DeepSeek-R1 and Qwen3) that emphasize advanced logical reasoning and test-time scaling over pure pattern matching \cite{deepseek2024v3, qwen32025}.

Concurrently, a pivotal paradigm shift has occurred: moving from static sequence-to-sequence generation toward interactive, tool-augmented autonomous agents. This transition was initially catalyzed by models pre-trained on vast codebases, such as Codex \cite{chen2021codex}, which demonstrated that LLMs could translate natural language into actionable programmatic commands. Building upon this, frameworks like ReAct \cite{yao2023react} and Toolformer \cite{schick2023toolformer} equipped models with the ability to synergize step-by-step reasoning with external API invocations. As reasoning engines matured, the focus shifted toward multi-agent collaboration and fully autonomous software engineering capabilities. Recent advancements, ranging from open-source frameworks like AutoGen and SWE-agent \cite{wu2023autogen, yang2024sweagent} to highly integrated proprietary systems like Claude Code \cite{anthropic2025claudecode}, showcase agents capable of navigating complex file systems, executing code iteratively, and autonomously self-correcting. Collectively, these advances established LLM agents as a viable substrate for complex, interactive applications — among which personalized education has emerged as a particularly compelling domain.

\subsection{Foundation Models for Education}
The integration of generative AI into educational technology represents a profound shift from static, rule-based Intelligent Tutoring Systems (ITS) to dynamic, generative frameworks. Early approaches primarily relied on rigid, operation-based formalisms to parse and solve domain-specific tasks, such as mathematical word problems \cite{amini2019mathqa}. However, the advent of Large Language Models has fundamentally upgraded this paradigm. Comprehensive surveys now outline the transition toward utilizing foundation models as core cognitive engines for personalized education, emphasizing their capacity to dynamically tailor feedback, simulate roleplay, and adapt to individual learning styles \cite{llmeducation2026}.

To bridge the gap between general-purpose text generation and strict pedagogical utility, recent advancements have focused on developing specialized, agentic tutoring frameworks. For instance, models such as MathCoder \cite{wang2024mathcoder} interleave natural language reasoning with verifiable code execution, ensuring that students receive accurate, step-by-step guidance rather than hallucinated solutions. Furthermore, the integration of LLMs with Knowledge Tracing (KT) paradigms enables these agents to actively model a student's evolving cognitive state. By maintaining dynamic representations of a learner's historical performance, state-of-the-art educational agents can autonomously adapt their pedagogical scaffolding—seamlessly shifting between direct hints and deeper Socratic questioning based on real-time comprehension \cite{macina2023mathdial}. Nevertheless, the practical deployment of such educational agents at scale remains constrained by two orthogonal challenges that have received comparatively less attention: the computational cost of running these models on accessible hardware, and their tendency to misalign with localized curricula — issues we now examine in turn.

\subsection{Post-Training Quantization for Efficient Deployment}
\label{subsec:quantization}

A fundamental prerequisite for deploying state-of-the-art LLMs on consumer-grade hardware is reducing the static memory footprint of model weights. Post-training quantization (PTQ) has emerged as the dominant solution, compressing pre-trained weights into low-bit representations without expensive re-training. GPTQ~\cite{frantar2023gptq} pioneered accurate one-shot quantization for generative transformers, employing an approximate second-order procedure to minimize layer-wise reconstruction error at 3 or 4 bits. AWQ~\cite{lin2024awq} further observes that weight importance is highly non-uniform: a small fraction of salient weights, identified through activation magnitudes, disproportionately influences output quality, and selectively preserving them throughper-channel scaling achieves superior accuracy under hardware-friendly uniform bit-widths. However, quantization addresses only the static weight footprint; the dynamic memory growth of the KV cache during long-context inference remains entirely unresolved, motivating the complementary techniques surveyed next.

\subsection{Knowledge Distillation for Compact Educational Models}
\label{subsec:distillation}

An alternative compression paradigm transfers the capabilities of large 
teacher models into smaller student architectures through knowledge 
distillation (KD). Traditional KD~\cite{hinton2015distilling} matches 
soft probability distributions, but often induces hallucinations in 
generative tasks due to the mode-covering behavior of forward 
Kullback-Leibler divergence. MiniLLM~\cite{gu2024minillm} mitigates 
this through reverse KLD with policy gradient optimization, forcing 
students to concentrate probability mass on high-quality generation 
regions. Beyond probability matching, Hsieh et al.~\cite{hsieh2023distilling} 
introduced step-by-step distillation that extracts intermediate reasoning 
trajectories from proprietary teachers, enabling compact students to 
internalize multi-step problem-solving heuristics --- a capability 
particularly valuable for pedagogical scaffolding in educational tutoring.

Nevertheless, distillation pipelines remain computationally intensive 
at training time and produce student models with static knowledge that 
cannot adapt to evolving curricula without repeated re-distillation, 
limiting their applicability for institutions facing both hardware and 
content-update constraints.

\subsection{Long-Context Inference Optimization}
\label{subsec:long-context}

Due to the quadratic computational complexity of standard attention, 
Transformer-based LLMs historically operate within limited pre-training 
context windows. Extending this window has been pursued through three 
complementary directions. First, modifications to positional encoding 
via interpolation and rotary embeddings enable zero-shot length 
generalization without re-training~\cite{chen2023extending, su2024roformer, 
ntk2023web, peng2024yarn, giraffe2023}. Second, long-context post-training 
explicitly adapts model weights to extended sequences, often supported 
by sequence parallelism techniques~\cite{gemini152024, sparseattention2024, 
du2024chatglm, shoeybi2020megatron}. Third, the integration of specialized 
memory modules enables processing of effectively unbounded 
sequences~\cite{dai2019transformerxl, rae2020compressive, bulatov2022recurrent, 
munkhdalai2024leave}.

However, even with these extensions, long-context inference suffers 
from prohibitive latency during the prefilling stage. System-level 
optimizations such as IO-aware FlashAttention and PagedAttention 
mitigate GPU memory bottlenecks but do not reduce the fundamental 
computational complexity of attention~\cite{dao2022flashattention, 
kwon2023efficient}. To process million-token contexts efficiently, 
a growing body of work exploits attention sparsity through 
input-independent patterns such as localized windows~\cite{wang2019multipassage, 
untieknots2025} and attention sinks that preserve critical initial 
tokens~\cite{han2024lminfinite, xiao2024efficient}. While effective, 
these heuristic patterns can permanently discard informative tokens, 
motivating dynamic KV cache selection methods that adaptively retain 
relevant entries~\cite{zhang2024h2o, li2024snapkv}. State-of-the-art 
frameworks further refine selection granularity, progressing from 
block-level memory units~\cite{xiao2024infllm}, through query-aware 
page-level sparsity~\cite{tang2024quest} and dynamic context 
selection~\cite{omnikv2025, jiang2024minference}, to fine-grained 
token-level selection that bypasses block-level inefficiencies 
entirely~\cite{tokenselect2025}.

While token-level methods such as TokenSelect substantially reduce 
per-step attention cost, the interaction between selective sparse 
attention and chunked prefill at very long contexts remains an active 
area of investigation.

\subsection{LLMs on Low-Resource Languages}
While modern LLMs exhibit remarkable zero-shot and few-shot reasoning capabilities in high-resource languages like English, their performance often degrades significantly when applied to low-resource languages (LRLs) such as Vietnamese. This degradation is primarily attributed to a severe underrepresentation of LRL tokens in the pre-training corpora, which limits the models' native linguistic alignment. 

To circumvent this linguistic bottleneck without resorting to computationally prohibitive pre-training from scratch, recent literature has explored several efficient adaptation strategies. At the inference level, Nguyen et al. \cite{democratizing2024acl} propose leveraging the models' dominant English capabilities through linguistically-diverse prompts. This approach acts as a cross-lingual bridge, effectively anchoring LRL inputs to the models' highly structured English latent space to improve zero-shot generalization. Complementarily, Cahyawijaya et al. \cite{cahyawijaya2024llms} demonstrate that LLMs possess latent cross-lingual capabilities that can be explicitly activated via In-Context Learning (ICL). Their empirical findings show that providing a limited number of task-specific demonstrations in the target language allows the model to dynamically adapt its reasoning pathways to the LRL without requiring any parameter updates. 

For applications requiring deeper structural alignment, prompt engineering alone is often insufficient. To address this, Nag et al. \cite{efficientcpt2025naacl} introduce efficient continual pre-training (CPT) techniques tailored for LRLs. Their methodology natively injects new linguistic knowledge into the foundation model's weights while actively mitigating the catastrophic forgetting of general reasoning capabilities, all under strict computational constraints. Collectively, these methodologies highlight that effective LRL 
adaptation can be achieved without invasive parameter updates, 
particularly when combined with carefully curated in-context signals.

\subsection{Hallucination in Localized and Cultural Contexts}
While modern LLMs demonstrate strong generalization capabilities, their deployment in region-specific applications is frequently hindered by factual hallucinations. Ji et al. \cite{ji2023survey} identify that hallucinations fundamentally stem from data divergence, where models learn spurious correlations from imbalanced pre-training corpora. This issue is severely amplified by the Western-centric nature of foundational models. As Huang et al. \cite{huang2023survey} highlight, when LLMs process queries outside their dominant data distribution, they often generate confident confabulations rather than acknowledging knowledge boundaries. Empirical evaluations by Bang et al. \cite{bang2023multitask} confirm this limitation, demonstrating that LLMs exhibit significantly higher hallucination rates and degraded reasoning when applied to low-resource languages and non-Western cultural contexts. In educational settings, relying on such ungrounded models introduces critical pedagogical risks. These findings establish that scaling pre-training alone cannot resolve 
the localization gap, particularly in educational deployments where 
factual accuracy on regional content is non-negotiable. This persistent 
gap motivates runtime grounding mechanisms capable of injecting 
localized knowledge without the prohibitive cost of continual 
re-training.

\subsection{Agentic Systems and Self-Evolving Context}
\label{subsec:agentic}
The rapid evolution of LLMs has introduced a paradox in knowledge 
management: while models store vast generalized knowledge within their 
parametric memory, their ability to seamlessly integrate deeply 
specialized domain heuristics remains limited. Traditional fine-tuning 
is rigid, computationally expensive, and susceptible to catastrophic 
forgetting. Conversely, current prompt optimization methods often 
suffer from brevity bias --- the tendency to compress prompts 
into short, generic instructions that strip away essential technical 
nuances required for high-precision tasks~\cite{promptalchemist2025}. 
Furthermore, attempts at monolithic rewriting in long-context windows 
frequently trigger context collapse, wherein iterative LLM 
rewrites progressively erode accumulated domain knowledge into 
uninformative summaries.

To overcome these barriers and build agentic systems that self-improve without weight updates, recent literature explores natural language feedback and verbal reinforcement learning \cite{krause2019dynamic, shinn2023reflexion, nature2025optimizing}. Techniques leveraging reflective prompt evolution \cite{gepa2025} and agentic context engineering \cite{agenticcontext2026} effectively transform the static context into a living, continuously updated "playbook." These methodologies enable LLM agents to adapt and fine-tune their behavior exclusively through dynamic context memory, preserving granular domain insights without the massive overhead of standard fine-tuning~\cite{memento_no_date}.

\subsection{Discussion and Positioning}
\label{subsec:rw-positioning}

Table~\ref{tab:positioning} summarizes where DeepEdu sits relative to the
four directions above. Each addresses one facet of on-premise educational
deployment but leaves the others open: quantization shrinks static weights
yet does nothing for the dynamic long-context cost; selective sparse
attention accelerates long contexts but neither self-improves nor
localizes; agentic context engineering self-improves without weight
updates but is not built for million-token efficiency or a specific
curriculum; and low-resource-language adaptation localizes, but through
weight updates that are costly to maintain as curricula evolve. DeepEdu is,
to our knowledge, the first framework to occupy their intersection: it
couples a cluster-level long-context engine with a self-improving,
curriculum-grounded agentic layer, so that a single on-premise system is
efficient, self-improving, and localizable at once.

\begin{table}[t]
\centering
\caption{Positioning of DeepEdu against representative prior directions.
\checkmark: directly addressed; (\checkmark): partial; --: not addressed.
Columns: on-premise/data-sovereignty fit, long-context efficiency,
self-improvement without fine-tuning, and localization.}
\label{tab:positioning}
\footnotesize
\setlength{\tabcolsep}{4pt}
\begin{tabular}{lcccc}
\toprule
\textbf{Direction} & \textbf{On-prem.} & \textbf{Long-ctx} & \textbf{Self-impr.} & \textbf{Local.} \\
\midrule
Quantization~\cite{lin2024awq, frantar2023gptq}                  & \checkmark & -- & -- & -- \\
Sparse long-context~\cite{tokenselect2025, xiao2024infllm}       & -- & \checkmark & -- & -- \\
Agentic context~\cite{agenticcontext2026, shinn2023reflexion}    & (\checkmark) & -- & \checkmark & -- \\
LRL adaptation~\cite{efficientcpt2025naacl, cahyawijaya2024llms} & -- & -- & -- & \checkmark \\
\midrule
\textbf{DeepEdu (ours)}                                          & \checkmark & \checkmark & \checkmark & \checkmark \\
\bottomrule
\end{tabular}
\end{table}

\section{Preliminaries}
\label{sec:preliminaries}

This section establishes the notation and the two paradigms our framework
builds on. We first summarize Transformer inference and Multi-Head
Attention (§\ref{subsec:llm-inference}) and formalize long-context
inference as a Selective Sparse Attention problem
(§\ref{subsec:problem-formulation})---the basis of our inference
engine---then define Agentic Context Engineering
(§\ref{subsec:prelim-ace}), the paradigm underlying our self-improving
agentic layer.

\subsection{LLM Inference and Multi-Head Attention}
\label{subsec:llm-inference}

\paragraph{Notation} 
Throughout this paper, we adopt the following unified notation. Let $d$ 
denote the model hidden dimension, $H$ the number of attention heads, 
and $d_h = d/H$ the per-head dimension. For a transformer layer processing 
$n$ input tokens, we denote its hidden states by 
$\mathbf{X} \in \mathbb{R}^{n \times d}$. The KV Cache accumulated from 
previously processed tokens has length $N$, and we use $h \in \{1, \ldots, H\}$ 
to index attention heads.

\paragraph{Multi-head attention}
Given hidden states $\mathbf{X}$, each head $h$ projects $\mathbf{X}$ into
per-head queries, keys, and values
$\mathbf{Q}^{(h)}, \mathbf{K}^{(h)}, \mathbf{V}^{(h)} \in \mathbb{R}^{n \times d_h}$
and computes scaled dot-product attention (SDPA)~\cite{vaswani2017attention}:
\begin{equation}
    \mathbf{O}^{(h)} = \text{softmax}\!\left(\frac{\mathbf{Q}^{(h)} (\mathbf{K}^{(h)})^\top}{\sqrt{d_h}}\right) \mathbf{V}^{(h)},
    \label{eq:sdpa}
\end{equation}
whose per-head outputs are concatenated and projected back to dimension
$d$. The inner softmax in Equation~\eqref{eq:sdpa} has quadratic
complexity $\mathcal{O}(n^2 d_h)$ in sequence length---the principal
bottleneck of long-context inference.

\paragraph{Prefill, decode, and the KV Cache}
Inference has two stages. \emph{Prefill} processes the whole prompt of
length $n_{in}$ in one pass and persists the per-head keys and values as a
KV Cache
$\mathbf{K}^{(h)}_{\text{cache}}, \mathbf{V}^{(h)}_{\text{cache}} \in \mathbb{R}^{N \times d_h}$
(initially $N = n_{in}$); it is compute-bound and dominates latency for
long inputs. \emph{Decode} then generates one token per pass, appending
its key and value to the cache so that $N$ grows by one each step. Each
decode step is light, but the KV Cache footprint grows linearly with
context, straining consumer-grade hardware.

\paragraph{Two-level chunking for long-context prefill}
To bound peak memory consumption during prefill and enable processing 
of inputs that exceed available VRAM, modern inference frameworks such 
as vLLM~\cite{kwon2023efficient} and SGLang adopt chunked prefill: 
the input sequence is partitioned into consecutive outer chunks, each 
of length at most $L_{\text{outer}}$ tokens (typically 
$L_{\text{outer}} = 8192$). Each outer chunk is processed in a separate 
forward pass; the last outer chunk may be shorter when $n_{in}$ is not 
a multiple of $L_{\text{outer}}$. This outer chunking caps the peak 
activation memory of any single forward pass but does not, by itself, 
reduce the cost of attending to the growing KV Cache: a naive 
implementation still incurs $\mathcal{O}(\ell \cdot N \cdot d_h)$ 
attention cost per outer chunk under Equation~\eqref{eq:sdpa}, where 
$\ell \leq L_{\text{outer}}$ denotes the actual length of the current 
outer chunk and $N$ is the cumulative KV Cache length.

To further reduce per-step attention cost, recent selective sparse 
attention methods such as TokenSelect~\cite{tokenselect2025} introduce 
a second, finer-grained level of partitioning \textit{within} each 
outer chunk. Specifically, the queries of an outer chunk of length 
$\ell$ are subdivided into $M = \lceil \ell / L \rceil$ sub-chunks, 
denoted $\{\mathbf{X}_1, \mathbf{X}_2, \ldots, \mathbf{X}_M\}$. Each 
sub-chunk has length $L$ (typically $L = 512$), except possibly the 
last, which has length $\ell - (M-1)L \leq L$. For each sub-chunk, the 
method invokes a query-aware selection function to identify a small 
subset of critical KV Cache entries, then computes sparse attention 
restricted to that subset. Crucially, this selection-then-attention 
procedure is applied sequentially to each sub-chunk within 
the outer chunk: $M$ independent retrieval calls followed by $M$ 
independent attention calls per outer chunk. While this scheme 
effectively bypasses the quadratic attention cost, the sequential 
nature of the inner loop exposes a substantial opportunity for 
amortization when consecutive sub-chunks share similar query 
representations --- a redundancy our work directly targets.

\subsection{Selective Sparse Attention}
\label{subsec:problem-formulation}

The quadratic cost of SDPA in Equation~\eqref{eq:sdpa}, combined with
the abundant sparsity empirically observed in LLM attention
maps~\cite{xiao2024infllm, tang2024quest, tokenselect2025}, motivates
the selective sparse attention (SSA) paradigm: rather than attending to
the entire KV Cache, only a small subset of critical tokens is
dynamically selected for each query. Our inference engine builds directly
on this paradigm; we restate it below in our notation, following the
formulation of TokenSelect~\cite{tokenselect2025}.

\begin{definition}[Selective Sparse Attention]
\label{def:ssa}
Consider a transformer layer at any inference step, with per-head 
queries $\mathbf{Q}^{(h)} \in \mathbb{R}^{C \times d_h}$ for $C$ current 
input tokens ($C = 1$ during decoding, $C \leq L$ for one sub-chunk during prefill where $L$ is the sub-chunk size defined in §\ref{subsec:llm-inference}), 
and a KV Cache 
$\mathbf{K}^{(h)}_{\text{cache}}, \mathbf{V}^{(h)}_{\text{cache}} \in \mathbb{R}^{N \times d_h}$. 
The full attention output following Equation~\eqref{eq:sdpa} is:
\begin{equation}
    \mathbf{O}^{(h)} = \text{softmax}\!\left(\frac{\mathbf{Q}^{(h)} (\mathbf{K}^{(h)}_{\text{cache}})^\top}{\sqrt{d_h}}\right) \mathbf{V}^{(h)}_{\text{cache}}.
    \label{eq:full-attention}
\end{equation}
Building on Equation~\eqref{eq:full-attention}, selective sparse attention
restricts the softmax to a subset of $k \ll N$ cache entries, yielding the
sparse approximation $\hat{\mathbf{O}}^{(h)}$:
\begin{equation}
    \hat{\mathbf{O}}^{(h)} = \text{softmax}\!\left(\frac{\mathbf{Q}^{(h)} (\mathbf{K}^{(h)}_{\text{select}})^\top}{\sqrt{d_h}}\right) \mathbf{V}^{(h)}_{\text{select}},
    \label{eq:sparse-attention}
\end{equation}
The two formulations differ only in \emph{which} cache rows enter the
softmax: the selected matrices
$\mathbf{K}^{(h)}_{\text{select}}, \mathbf{V}^{(h)}_{\text{select}} \in \mathbb{R}^{k \times d_h}$
index $\mathbf{K}^{(h)}_{\text{cache}}$ and
$\mathbf{V}^{(h)}_{\text{cache}}$ at the positions $\mathcal{I}$ chosen by
a \textit{selection function} $\mathcal{S}$:
\begin{equation}
    \mathcal{I} = \mathcal{S}\!\left(\mathbf{Q}, \mathbf{K}_{\text{cache}}\right), \quad \mathcal{I} \subseteq \{1, \ldots, N\}, \quad |\mathcal{I}| = k.
    \label{eq:selection}
\end{equation}
Here $\mathbf{Q}$ and $\mathbf{K}_{\text{cache}}$ denote the multi-head 
queries and cached keys aggregated across all heads, allowing $\mathcal{S}$ 
to leverage cross-head information when scoring token criticality. The 
design objective is to construct $\mathcal{S}$ such that the approximation 
gap $\big\|\mathbf{O}^{(h)} - \hat{\mathbf{O}}^{(h)}\big\|_2^2$ is 
minimized while keeping $k$ within the hardware budget.
\end{definition}

Selection functions $\mathcal{S}$ range from fixed patterns (attention
sinks and recent tokens), through query-independent scoring, to query-aware
token-level selection (surveyed in §\ref{subsec:long-context}). Our engine
builds on the query-aware token-level scheme of
TokenSelect~\cite{tokenselect2025}, which conditions $\mathcal{S}$ on the
current query $\mathbf{Q}$ for high approximation quality, at the cost of a
retrieval call per step.

\paragraph{Sub-chunk redundancy in selective sparse attention}
Together, Equations~\eqref{eq:sparse-attention}--\eqref{eq:selection}
define a single \emph{selection-then-attention} step, the unit our
inference engine reorganizes. A critical observation underexplored by
prior work is that, under the two-level chunking scheme of
§\ref{subsec:llm-inference}, this step---and in particular the costly
selection function $\mathcal{S}$---is invoked independently for each of
the $M$ inner sub-chunks
$\{\mathbf{X}_1, \ldots, \mathbf{X}_M\}$ that constitute a single outer 
prefill chunk. Each invocation triggers a full pass over the growing KV 
Cache to score and retrieve the top-$k$ critical tokens, followed by an 
attention call on the selected subset. When consecutive sub-chunks 
exhibit highly similar query representations --- as is often the case 
within a coherent passage of text --- a substantial fraction of their 
retrieved index sets overlap, yet the retrieval cost is incurred in 
full for each sub-chunk independently. This per-sub-chunk invocation 
pattern induces redundant retrieval computation that grows linearly 
with $M$, a phenomenon we empirically analyze in 
Section~\ref{sec:motivations} and address through our proposed 
\textit{cluster-level} selection paradigm in 
Section~\ref{sec:methodology}, which groups similar consecutive 
sub-chunks before invoking $\mathcal{S}$ a single time per cluster.

\subsection{Agentic Context Engineering}
\label{subsec:prelim-ace}

While selective sparse attention addresses the computational bottleneck,
the semantic bottleneck---grounding the model in localized, trustworthy
knowledge---calls for a complementary paradigm. Our agentic layer builds
on Agentic Context Engineering (ACE)~\cite{agenticcontext2026}, which
adapts an LLM agent by editing its context rather than its weights. We
restate it in the notation used later.

\begin{definition}[Agentic Context Engineering]
\label{def:ace}
Let a backbone LLM with fixed parameters $\phi$ process a stream of
interactions $\mathcal{D}=\{(q_t,c_t,y_t)\}_{t=1}^{T}$, where $q_t$ is a
query, $c_t$ its available context, and $y_t$ a feedback signal (e.g., a
reference answer or a deterministic environment outcome). ACE maintains a
structured \emph{playbook} $\mathcal{P}_t$ of itemized natural-language
entries and, after each interaction, applies a validated set of edits
$\Delta_t$ distilled from the resulting trajectory:
\begin{equation}
    \mathcal{P}_{t+1}=\operatorname{Apply}(\mathcal{P}_t,\Delta_t),
    \qquad \phi_{t+1}=\phi_t.
    \label{eq:ace-context-update}
\end{equation}
Adaptation therefore occurs entirely in inspectable context, leaving the
weights $\phi$ unchanged.
\end{definition}

Equation~\eqref{eq:ace-context-update} leaves open two design choices that
decide whether such an agent actually improves: \emph{how} the edits
$\Delta_t$ are produced, and \emph{which} entries condition each
interaction as the playbook grows. Sections~\ref{sec:motivations}
and~\ref{sec:methodology} resolve these in turn---the former through
empirical observations, the latter through the concrete
Generator--Reflector--Curator design of our agentic layer.

\section{Motivations and Observations}
\label{sec:motivations}

The design of SCALE is grounded in five empirical observations, three
concerning the long-context engine and two concerning the agentic layer.
First, consecutive query sub-chunks within an outer chunk exhibit high
mutual similarity across the depth of the model
(§\ref{subsec:obs-similarity}). Second, leveraging this similarity
through a tunable threshold $\theta$ enables aggressive cluster
formation, with a controllable trade-off between amortization and
selection granularity (§\ref{subsec:obs-cluster}). Third, the top-$k$
indices retrieved by a single cluster-level query closely match those
that would be retrieved per-sub-chunk, confirming that cluster-level
amortization preserves selection fidelity
(§\ref{subsec:obs-topk-overlap}). Fourth, as the Playbook grows,
supplying each Generator call with only the query-relevant rules is both
faster and more accurate than passing the entire Playbook, and a
length-matched random control confirms the gain comes from query--rule
matching rather than from shorter prompts alone
(§\ref{subsec:obs-rae}). Fifth, auditing the
agent's incorrect trajectories reveals that its failures are not
isolated accidents but recur across superficially different tasks, so
the same underlying weaknesses resurface until they are explicitly
addressed (§\ref{subsec:obs-agentic}). Together, these observations directly
motivate the cluster-level sparse selection of our long-context engine
(§\ref{subsec:long-context-engine}) and the playbook-evolving design of our
agentic layer (§\ref{subsec:agentic-engine}).


\subsection{Consecutive Sub-Chunk Similarity within Outer Chunks}
\label{subsec:obs-similarity}

Under the chunked prefill scheme of §\ref{subsec:llm-inference}, each 
outer chunk of length $\ell$ is subdivided into 
$M = \lceil \ell / L \rceil$ sub-chunks 
$\{\mathbf{X}_1, \ldots, \mathbf{X}_M\}$ with $L = 512$. We empirically 
investigate whether the query representations of these sub-chunks 
remain similar enough across $i = 1, \ldots, M-1$ to justify 
amortizing retrieval across them. Following 
TokenSelect~\cite{tokenselect2025}, we summarize each sub-chunk 
$\mathbf{X}_i$ by its mean-pooled query representation:
\begin{equation}
    \mathbf{v}_i = \frac{1}{|\mathbf{X}_i|} \sum_{t=1}^{|\mathbf{X}_i|} \mathbf{q}_{i,t},
    \label{eq:chunk-fingerprint}
\end{equation}
where $\mathbf{q}_{i,t} \in \mathbb{R}^{d}$ denotes the multi-head query
of the $t$-th token in sub-chunk $\mathbf{X}_i$, so that
$\mathbf{v}_i \in \mathbb{R}^{d}$; we write
$\mathbf{v}_i^{(h)} \in \mathbb{R}^{d_h}$ for its head-$h$ slice. The pairwise similarity between consecutive 
sub-chunks is measured by cosine similarity 
$\cos(\mathbf{v}_i, \mathbf{v}_{i+1})$.

To probe this, we profile prefill over six representative long-context
tasks from InfiniteBench~\cite{zhang_etal_2024_bench} (Code.D, R.KV,
En.Dia, Math.F, R.Num, and R.PK) using
Qwen2-7B-Instruct~\cite{yang2024qwen2} as the backbone
($L_{\text{outer}} = 8192$, $L = 512$), recording
$\cos(\mathbf{v}_i, \mathbf{v}_{i+1})$ for every consecutive sub-chunk
pair across five layers that span the depth of the model: shallow
(Layers~0 and~6), middle (Layers~13 and~20), and deep (Layer~27).

Table~\ref{tab:sim-distribution} summarizes the resulting distribution,
and the picture is unambiguous: consecutive sub-chunks remain highly
similar throughout the network. The median cosine similarity stays
above $0.92$ at every layer---exceeding $0.94$ at the shallow and deep
layers---and even the $10$th percentile never falls below $0.86$,
indicating that high similarity is the prevailing regime rather than
an artifact of a few dominant pairs. Mass concentrates at the upper
end of the distribution and thins only at the extreme lower tail, where
a small number of transitions drop sharply, down to $0.33$ at Layer~13.
We attribute these rare drops to genuine semantic shifts---such as
boundaries between distinct documents or context segments---rather than
pervasive noise, an interpretation corroborated in
§\ref{subsec:obs-cluster}, where our anchor-based clustering closes a
cluster precisely at such transitions.

\begin{table}[t]
\centering
\caption{Distribution of consecutive sub-chunk cosine similarity 
$\cos(\mathbf{v}_i, \mathbf{v}_{i+1})$, aggregated across six 
long-context tasks (Qwen2-7B-Instruct, $L = 512$). Consecutive
sub-chunks share nearly identical query representations at every layer,
with the low-similarity tail confined to rare semantic boundaries.}
\label{tab:sim-distribution}
\small
\begin{tabular}{lcccccc}
\toprule
Layer & Mean & Min & Max & P10 & Median & P90 \\
\midrule
Layer 0  & 0.938 & 0.442 & 1.000 & 0.887 & 0.943 & 0.982 \\
Layer 6  & 0.956 & 0.797 & 1.000 & 0.934 & 0.958 & 0.977 \\
Layer 13 & 0.925 & 0.329 & 0.995 & 0.876 & 0.934 & 0.969 \\
Layer 20 & 0.918 & 0.580 & 0.996 & 0.868 & 0.924 & 0.963 \\
Layer 27 & 0.943 & 0.434 & 1.000 & 0.870 & 0.959 & 0.987 \\
\bottomrule
\end{tabular}
\end{table}

This concentration of mass at high similarity suggests that treating
each sub-chunk as an independent retrieval target is largely
redundant: the query representations driving these retrievals are
nearly identical for most consecutive pairs. Two design questions 
follow: \textit{how aggressively} can we group similar sub-chunks into 
clusters, and \textit{at what cost} to selection fidelity? We address 
these questions in turn in the next two subsections.


\subsection{Cluster Formation Under Tunable Similarity Threshold}
\label{subsec:obs-cluster}

The high mutual similarity established in §\ref{subsec:obs-similarity} 
suggests that consecutive sub-chunks can be grouped into clusters 
whenever their pairwise similarity exceeds a threshold $\theta$. We 
empirically characterize how the choice of $\theta$ controls the 
aggressiveness of cluster formation across the depth of the model, 
exposing a tunable accuracy--efficiency trade-off.

To this end, for each outer chunk processed during prefill we apply the
anchor-based clustering described in §\ref{sec:methodology}---consecutive
sub-chunks join a cluster as long as their similarity to the cluster
anchor exceeds $\theta$, up to a maximum cluster size of
$L_{\max} = 8 \times L = 4096$ tokens---and record the resulting cluster
sizes as $\theta$ is swept over $\{0.95, 0.97, 0.99, 0.999\}$.

Table~\ref{tab:cluster-sizes} reports the mean cluster size at each
threshold. The mechanism is markedly \textit{depth-dependent}: shallow
layers (Layers~0 and~6) and the final layer (Layer~27) sustain large
clusters across a wide range of $\theta$, whereas the middle layers
(Layers~13 and~20) fragment far more readily as $\theta$ tightens---at
$\theta = 0.99$, for instance, the mean cluster at Layer~20 collapses to
$1628$ tokens while Layer~0 remains near-maximal at $4013$. This
asymmetry is consistent with prior findings that mid-depth layers in
autoregressive LLMs encode richer semantic representations than either
shallow or final layers~\cite{skean2025layerbylayer}: the more
discriminative the representation, the more readily our mechanism
detects a genuine semantic transition and closes the cluster.

\begin{table}[t]
\centering
\caption{Mean cluster size (in tokens) at each threshold $\theta$. 
Higher $\theta$ produces smaller clusters as the mechanism increasingly 
detects semantic boundaries; the effect is most pronounced at middle 
layers where semantic discrimination is highest. Maximum cluster size 
is $L_{\max} = 4096$.}
\label{tab:cluster-sizes}
\small
\begin{tabular}{lcccc}
\toprule
Layer & $\theta = 0.95$ & $\theta = 0.97$ & $\theta = 0.99$ & $\theta = 0.999$ \\
\midrule
Layer 0  & 4025 & 4025 & 4013 & 3446 \\
Layer 6  & 4025 & 4021 & 3904 & 1616 \\
Layer 13 & 3988 & 3699 & 1947 & \phantom{0}829 \\
Layer 20 & 3964 & 3260 & 1628 & \phantom{0}694 \\
Layer 27 & 4017 & 3992 & 3695 & 1509 \\
\bottomrule
\end{tabular}
\end{table}

Zooming in on Layer~20, the most threshold-sensitive layer,
Table~\ref{tab:cluster-merge-rate} reports the fraction of clustering
decisions that yield cluster sizes of $1, 2, \ldots, 8$ blocks. At the
default operating point $\theta = 0.95$, more than
$94\%$ of clustering decisions reach the maximum cluster size of $8$ 
blocks, reducing the number of retrieval invocations of $\mathcal{S}$ 
per outer chunk by approximately $8\times$ relative to the 
per-sub-chunk baseline. As $\theta$ tightens to $0.999$, the mechanism 
gradually reverts to fine-grained per-sub-chunk selection, with 
single-block decisions dominating ($82.6\%$). This continuous spectrum 
exposes $\theta$ as a principled tuning knob: lower $\theta$ maximizes 
amortization, higher $\theta$ maximizes selection granularity.

\begin{table}[t]
\centering
\caption{Distribution of cluster sizes (in $L$-blocks) at Layer 20 
under varying threshold $\theta$. At $\theta = 0.95$, $94\%$ of 
clusters reach the maximum size of $8$ blocks; at $\theta = 0.999$, 
the mechanism reverts to single-block selection $83\%$ of the time.}
\label{tab:cluster-merge-rate}
\small
\begin{tabular}{lcccc}
\toprule
Cluster size & $\theta\!=\!0.95$ & $\theta\!=\!0.97$ & $\theta\!=\!0.99$ & $\theta\!=\!0.999$ \\
\midrule
1 block        & \phantom{0}0.6\% & \phantom{0}2.4\% & 36.7\% & 82.6\% \\
2 blocks       & \phantom{0}0.9\% & \phantom{0}7.2\% & 19.2\% & \phantom{0}9.2\% \\
3 blocks       & \phantom{0}1.4\% & \phantom{0}8.2\% & 11.0\% & \phantom{0}3.7\% \\
4 blocks       & \phantom{0}1.5\% & \phantom{0}8.0\% & \phantom{0}8.2\% & \phantom{0}2.1\% \\
5--7 blocks    & \phantom{0}1.3\% & 12.4\% & \phantom{0}9.3\% & \phantom{0}2.0\% \\
8 blocks (max) & 94.2\% & 61.7\% & 15.5\% & \phantom{0}0.5\% \\
\bottomrule
\end{tabular}
\end{table}

These results establish that aggressive cluster formation is 
\textit{achievable} under realistic operating conditions: at 
$\theta = 0.95$, the dominant operating point of our mechanism, the 
overwhelming majority of clustering decisions yield maximum-size 
clusters. The remaining question is whether this aggressive 
amortization preserves the fidelity of the underlying selection 
function $\mathcal{S}$, which we examine next.


\subsection{Top-$k$ Overlap Under Cluster-Level Retrieval}
\label{subsec:obs-topk-overlap}

The first two observations established that consecutive sub-chunks 
are highly similar (§\ref{subsec:obs-similarity}) and that this 
similarity enables aggressive cluster formation 
(§\ref{subsec:obs-cluster}). We now verify the critical remaining 
question: when the cluster's anchor query is used in place of
per-sub-chunk queries, does the resulting top-$k$ retrieval recover the
same critical KV Cache entries as the per-sub-chunk baseline?

To answer this, for each cluster $\mathcal{C} = \{\mathbf{X}_i, \mathbf{X}_{i+1},
\ldots, \mathbf{X}_{i+|\mathcal{C}|-1}\}$ formed under threshold 
$\theta$, we perform two retrievals against the same KV Cache:
\begin{itemize}
    \item \textbf{Per-sub-chunk}: retrieve the top-$k$ index sets 
    $\mathcal{I}_i, \mathcal{I}_{i+1}, \ldots, \mathcal{I}_{i+|\mathcal{C}|-1}$ 
    independently using $\mathbf{v}_i, \mathbf{v}_{i+1}, \ldots$, then 
    take their union 
    $\mathcal{I}_{\text{union}} = \bigcup_{j} \mathcal{I}_{i+j}$.
    \item \textbf{Anchor (cluster-level)}: retrieve a single top-$k$
    index set $\mathcal{I}_{\mathcal{C}}$ using only the \emph{anchor
    fingerprint} $\mathbf{v}_i$---the query representation of the first
    sub-chunk, against which all other members were admitted---shared
    across the whole cluster, exactly as SCR does at inference time
    (§\ref{subsubsec:scr-retrieval}).
\end{itemize}
We measure the \textit{overlap ratio}:
\begin{equation}
    \rho(\mathcal{C}) = \frac{|\mathcal{I}_{\mathcal{C}} \cap \mathcal{I}_{\text{union}}|}{|\mathcal{I}_{\mathcal{C}}|},
    \label{eq:overlap-ratio}
\end{equation}
which quantifies the fraction of anchor-retrieved indices that would
also have been retrieved by independent per-sub-chunk retrievals. We
sweep $\theta \in \{0.95, 0.97, 0.99\}$, matching the operating points 
of §\ref{subsec:obs-cluster}, with $k = 8192$.

As reported in Table~\ref{tab:topk-overlap}, the overlap is
consistently high and rises monotonically with $\theta$: the median
$\rho$ already reaches $97.9\%$ at $\theta = 0.95$ and exceeds $99\%$
at $\theta = 0.99$, while the minimum stays above $88\%$ across all
measured clusters. In other words, throughout the operating range of
our mechanism a single anchor-driven retrieval recovers nearly all of
the critical tokens that per-sub-chunk retrieval would have selected
independently, rendering the $|\mathcal{C}|$ separate retrievals
largely redundant.

\begin{table}[t]
\centering
\caption{Top-$k$ overlap ratio $\rho$ (\%) between the anchor query and
per-sub-chunk retrieval at varying similarity thresholds $\theta$,
with $k = 8192$. Higher $\theta$ yields tighter overlap, confirming
that anchor-driven retrieval faithfully approximates per-sub-chunk
retrieval in the operating regime of our mechanism.}
\label{tab:topk-overlap}
\small
\begin{tabular}{lcccc}
\toprule
$\theta$ & Min & Median & Mean & Max \\
\midrule
$0.95$ & 88.2 & 97.9 & 96.8 & 100.0 \\
$0.97$ & 90.5 & 98.6 & 97.7 & 100.0 \\
$0.99$ & 93.1 & 99.3 & 98.6 & 100.0 \\
\bottomrule
\end{tabular}
\end{table}

The three observations together establish a chain of evidence for the
cluster-level paradigm. The high consecutive similarity 
(§\ref{subsec:obs-similarity}) creates the opportunity; the aggressive 
cluster formation at $\theta = 0.95$ (§\ref{subsec:obs-cluster}) 
demonstrates that the opportunity is realizable in practice; and the 
high top-$k$ overlap (§\ref{subsec:obs-topk-overlap}) confirms that 
amortization does not compromise selection fidelity. Our long-context 
inference engine (§\ref{sec:methodology}) operationalizes these 
insights through the Similarity Chunk Rolling mechanism, adopting 
$\theta = 0.95$ as the default operating point.

\subsection{Query-Aware Retrieval for Efficient Playbook Execution}
\label{subsec:obs-rae} 
As ACE accumulates more rules, supplying the entire Playbook to every
Generator call makes execution increasingly expensive and exposes the
model to guidance that is unrelated to the current task. The relevant
question is therefore not only whether a shorter prompt is faster, but
whether the rules retained in that prompt are the ones needed for the
current query.

We evaluate a fixed, fully-curated Formula Playbook in eval only mode (distinct from the un-adapted Origin Playbook of Section~\ref{subsec:exp-agentic}) with the same
Qwen3-4B-Instruct-2507~\cite{qwen32025} backbone, decoding configuration, and 200 test
examples in all conditions. The full-Playbook condition supplies every
entry to the Generator. Semantic RAE supplies the top $K=10$ entries
most relevant to the query and context. As a length-matched control,
random retrieval supplies ten randomly selected entries.

\begin{figure}[H]
\centering
\includegraphics[width=\columnwidth]{figures/time_rae.png}
\caption{Evaluation-only RAE study on Formula. All conditions use the
same fixed Playbook, Qwen3-4B-Instruct-2507, and 200 test examples.
Semantic top-$10$ retrieval is both more accurate and faster than the
full-Playbook condition; random top-$10$ is a length-matched control.}
\label{fig:obs-rae}
\end{figure}

Figure~\ref{fig:obs-rae} shows that semantic top-$10$ retrieval raises
accuracy from $74.5\%$ with the full Playbook to $79.5\%$, while
reducing mean TTFT from $0.392$ to $0.068$ seconds and mean end-to-end
generation time from $5.014$ to $2.319$ seconds. The random top-$10$
control retains the shorter context but reaches only $71.5\%$ accuracy.
Thus, the gain is not explained by truncating the Playbook alone: the
quality of query--rule matching is essential. Random retrieval lowers
TTFT and TPOT almost as much as semantic retrieval, yet it removes
useful guidance and performs worse than the full-Playbook condition.
Conversely, passing the entire Playbook preserves all potentially useful
rules but increases prefill cost and introduces distractors. These
results directly motivate Retrieval-Augmented Execution detailed in §\ref{subsec:agentic-engine}.

\subsection{Systematic and Recurring Agent Failures}
\label{subsec:obs-agentic}
An accumulated Playbook gives ACE persistent knowledge that can be
reused across tasks, but it does not guarantee that the agent's
underlying weaknesses have actually been resolved. When the base agent
fails a task for a systematic reason, the same weakness can resurface
on later, superficially different tasks. We therefore ask whether such
failures are isolated accidents or recurring, related events.


\begin{figure}[H]
\centering
\includegraphics[width=0.96\columnwidth]{figures/noise.png}
\caption{Recurring failure observations in the 500-task offline Formula
run, which produces 126 Ground-Truth-verified failures. Panel (a)
summarizes these verified failures, and panel (b) counts, for each
failure observation, how many later failure observations are
semantically related to it.}
\label{fig:noise}
\end{figure}


The finding is that a substantial fraction of tasks produces a
verified failure, and that these failures are often related across
distinct tasks. As shown in Figure~\ref{fig:noise}, in the 500-task
offline Formula run, 126 tasks produce a Ground-Truth-verified failure
($25.2\%$ of tasks). Of these failure observations, 91 ($72.2\%$) are related to at
least one later observation, while 55 are related to at least six later
observations. This pattern indicates that some underlying weaknesses can
recur even when their surface questions differ. 

The implication is that these recurring weaknesses should be exposed
proactively, rather than only when a future user query happens to
encounter them. A static Playbook cannot distinguish a robust rule from
one that merely appears plausible until it is tested against a boundary
case or a conflicting instruction, so a weakness that once caused a
failure can remain latent in the accumulated context and resurface on
later tasks until it is explicitly revised or removed. This finding
directly motivates the adversarial agent and failure-memory grounding in
our self-improvement loop detailed in §\ref{subsec:agentic-engine}.

\section{Methodology}
\label{sec:methodology}

This section presents SCALE, the framework at the technical core of
DeepEdu. Building on the empirical observations of §\ref{sec:motivations},
it comprises two synergistic components that jointly address the
computational and semantic challenges of deploying agentic LLMs in
resource-constrained educational environments;
Figure~\ref{fig:scale-overview} gives the end-to-end picture.
§\ref{subsec:long-context-engine} introduces \textit{Similarity Chunk
Rolling} (SCR), the long-context inference engine that amortizes selective
sparse attention (Definition~\ref{def:ssa}) from per-sub-chunk to
per-cluster granularity, while §\ref{subsec:agentic-engine} introduces the
self-improving agentic layer that instantiates the ACE paradigm
(Definition~\ref{def:ace}) to continuously curate a localized educational
playbook without expensive weight updates. The two are orthogonal yet
complementary: the engine keeps the long, playbook-conditioned contexts
curated by the agentic layer tractable under strict hardware budgets,
while the agentic layer ensures that efficient inference stays grounded in
locally accurate knowledge.

\begin{figure*}[t]
    \centering
    \includegraphics[width=0.9\textwidth]{figures/draft_scale.png}
    \caption{SCALE pipeline comparison. As the agentic layer continuously curates an evolving playbook
(10K\,$\to$\,200K tokens), the choice of inference engine determines deployment feasibility. Without
the long-context engine, prefill latency grows quadratically and accuracy degrades on noisy long
contexts, eventually triggering out-of-memory failures. SCALE's cluster-level sparse selection
maintains low latency and improves accuracy as the playbook accumulates richer localized knowledge.
Latency and accuracy numbers are illustrative; concrete benchmarks are reported in Section~\ref{sec:experiments}.}
    \label{fig:scale-overview}
\end{figure*}


\subsection{Efficient Long-Context Inference Engine}
\label{subsec:long-context-engine}

This subsection presents Similarity Chunk Rolling (SCR), a mechanism that
reorganizes the selection-then-attention step of
Equations~\eqref{eq:sparse-attention}--\eqref{eq:selection} so that it runs
\emph{once per cluster} of similar sub-chunks rather than once per
sub-chunk (§\ref{subsec:problem-formulation}). SCR dynamically groups
consecutive sub-chunks with high mutual similarity into clusters, invokes
the selection function $\mathcal{S}$ a single time per cluster, and
computes the sparse attention of Equation~\eqref{eq:sparse-attention} over
the entire cluster in one parallel forward pass.
Figure~\ref{fig:scr-pipeline} illustrates the overall procedure.

\begin{figure*}[t]
    \centering
    \includegraphics[width=0.9\textwidth]{figures/SCR_draft.png}
    \caption{Similarity Chunk Rolling (SCR) pipeline. Consecutive query
    sub-chunks with high mutual similarity are grouped into clusters via
    anchor-based formation; SCR issues a single top-$k$ retrieval per
    cluster and computes sparse attention over the whole cluster in one
    parallel forward pass, amortizing selection from per-sub-chunk to
    per-cluster granularity.}
    \label{fig:scr-pipeline}
\end{figure*}

\subsubsection{Cluster representation and similarity measure}
\label{subsubsec:scr-representation}

SCR represents each cluster through the sub-chunk fingerprints
$\mathbf{v}_i$ of Equation~\eqref{eq:chunk-fingerprint}. To make the
similarity checks in the clustering inner loop cheap, we pre-normalize
each fingerprint, $\tilde{\mathbf{v}}_i = \mathbf{v}_i / \|\mathbf{v}_i\|_2$,
so that cosine similarity reduces to a single inner product
$\cos(\mathbf{v}_i, \mathbf{v}_j) = \tilde{\mathbf{v}}_i^\top \tilde{\mathbf{v}}_j$,
implementable as one fused dot-product kernel with no repeated
L2-norm computation.

\subsubsection{Anchor-based cluster formation}
\label{subsubsec:scr-clustering}

A naive approach to grouping similar consecutive sub-chunks would 
compare each sub-chunk pairwise with its immediate predecessor: 
extend the current cluster whenever 
$\cos(\mathbf{v}_i, \mathbf{v}_{i+1}) \geq \theta$. However, this 
sliding-pair comparison is vulnerable to a phenomenon we 
call drift: small incremental changes between adjacent 
sub-chunks may individually pass the threshold, yet accumulate across 
the cluster, causing the final sub-chunk to diverge substantially 
from the cluster's starting point. The resulting cluster 
representation would no longer faithfully approximate any of its 
members.

To prevent this drift, we adopt an anchor-based procedure: the first
sub-chunk $\mathbf{X}_a$ of a cluster is fixed as the anchor, and each
subsequent $\mathbf{X}_j$ is admitted only while it stays similar to the
\emph{anchor}---not the drifting predecessor---and the cluster stays within
a per-call memory cap $L_{\max}$ (typically $L_{\max}=4096=8L$):
\begin{equation}
    \cos(\mathbf{v}_a, \mathbf{v}_j) \geq \theta \quad \text{and} \quad |\mathcal{C}_a| \cdot L < L_{\max}.
    \label{eq:scr-admission}
\end{equation}
The cluster closes at the first $\mathbf{X}_j$ violating either condition,
which then becomes the next anchor---so every member is similar to a
single fixed reference rather than a drifting predecessor, and no
attention call exceeds $L_{\max}$ tokens.
Algorithm~\ref{alg:scr-clustering} formalizes the full clustering 
procedure, which partitions the $M$ sub-chunks of an outer chunk 
into a set of disjoint clusters $\{\mathcal{C}_1, \mathcal{C}_2, \ldots\}$ 
in a single linear pass.

\begin{algorithm}[t]
\caption{Anchor-based clustering of sub-chunks.}
\label{alg:scr-clustering}
\begin{algorithmic}[1]
\Require Sub-chunk fingerprints $\{\tilde{\mathbf{v}}_1, \ldots, \tilde{\mathbf{v}}_M\}$; threshold $\theta$; max cluster size $L_{\max}$
\Ensure Clusters $\{\mathcal{C}_1, \mathcal{C}_2, \ldots\}$
\State $\mathcal{P} \gets \emptyset$, $a \gets 1$
\While{$a \leq M$}
    \State $\mathcal{C}_a \gets \{\mathbf{X}_a\}$, $j \gets a + 1$
    \While{$j \leq M$ and $\tilde{\mathbf{v}}_a^\top \tilde{\mathbf{v}}_j \geq \theta$ and $|\mathcal{C}_a| \cdot L < L_{\max}$}
        \State $\mathcal{C}_a \gets \mathcal{C}_a \cup \{\mathbf{X}_j\}$
        \State $j \gets j + 1$
    \EndWhile
    \State $\mathcal{P} \gets \mathcal{P} \cup \{\mathcal{C}_a\}$
    \State $a \gets j$
\EndWhile
\State \Return $\mathcal{P}$
\end{algorithmic}
\end{algorithm}

The output partition $\mathcal{P}$ contains $|\mathcal{P}| \leq M$ 
clusters; in the regime characterized by 
§\ref{subsec:obs-cluster}---where the majority of clustering 
decisions yield maximum-size clusters at $\theta = 0.95$---the 
reduction in retrieval invocations approaches the upper bound
$|\mathcal{P}| \approx M / (L_{\max}/L)$.

\paragraph{Cost model.} For an outer chunk of $M$ sub-chunks against a KV
Cache of length $N$, the per-sub-chunk baseline invokes the selection
function $\mathcal{S}$ once per sub-chunk, incurring $\mathcal{O}(MNd)$
scoring cost; SCR invokes it once per cluster, reducing this to
$\mathcal{O}(|\mathcal{P}|Nd)$. At the default operating point
($\theta\!=\!0.95$, $L_{\max}\!=\!8L$), where most clusters reach maximum
size, $|\mathcal{P}|\approx M/8$---an $8\times$ reduction in retrieval
scoring, which we measure empirically at $7.7\times$
(Figure~\ref{fig:prefill-breakdown}). The added bookkeeping---mean-pooling $M$
fingerprints and one linear clustering pass
(Algorithm~\ref{alg:scr-clustering})---is $\mathcal{O}(Md)$ and negligible
against a single retrieval's $\mathcal{O}(Nd)$ since $M\ll N$. The
sparse-attention FLOPs are unchanged: SCR attends to the same per-query
budget as the baseline but batches an entire cluster into one kernel
launch instead of $|\mathcal{C}_a|$ separate calls. Because the retrieval
term dominates and scales with $N$, the speedup grows with context
length.

\subsubsection{Cluster-level retrieval and parallel attention}
\label{subsubsec:scr-retrieval}

For each cluster $\mathcal{C}_a$ formed by
Algorithm~\ref{alg:scr-clustering}, we perform a \textit{single}
selection-then-attention step over the entire cluster. Rather than
aggregating its members, we reuse the cluster's \textit{anchor
fingerprint} $\mathbf{v}_a$---the query representation of the first
sub-chunk $\mathbf{X}_a$, against which every other member was admitted
under Equation~\eqref{eq:scr-admission}---directly as the cluster-level
query. The selection function $\mathcal{S}$ from
Equation~\eqref{eq:selection} is thus invoked once, with $\mathbf{v}_a$
acting as a single-row query:
\begin{equation}
    \mathcal{I}_{\mathcal{C}} = \mathcal{S}(\mathbf{v}_a, \mathbf{K}_{\text{cache}}),
    \label{eq:scr-cluster-selection}
\end{equation}
producing a single index set
$\mathcal{I}_{\mathcal{C}} \subseteq \{1, \ldots, N\}$ with
$|\mathcal{I}_{\mathcal{C}}| = k$ that is shared by all members of the
cluster. This single retrieval replaces the $|\mathcal{C}|$ independent
retrievals that the per-sub-chunk baseline would perform. Using the
anchor rather than a mean is deliberate: the admission criterion
Equation~\eqref{eq:scr-admission} guarantees
$\cos(\mathbf{v}_a, \mathbf{v}_j) \geq \theta$ for every member
$\mathbf{X}_j \in \mathcal{C}_a$, so all members share nearly collinear
query representations, and the token ranking induced by $\mathbf{v}_a$
closely matches the ranking each member would produce on its own. The
same threshold $\theta$ that forms the cluster therefore also controls
the retrieval approximation error---a fidelity we quantify empirically
in §\ref{subsec:obs-topk-overlap}.

\paragraph{Selection function instantiation} 
We instantiate $\mathcal{S}$ using TokenSelect's head soft-voting 
strategy~\cite{tokenselect2025}: per-head dot-product scores
$\mathbf{v}_a^{(h)} (\mathbf{K}^{(h)}_{\text{cache}})^\top$
are normalized through softmax independently, then summed across 
heads to produce a single criticality score per cached token, from 
which the top-$k$ indices are selected. The per-head softmax 
normalization ensures that no single head with disproportionately 
large attention logits dominates the selection.

\paragraph{Sparse attention over the cluster} 
Once $\mathcal{I}_{\mathcal{C}}$ is obtained, sparse attention is 
computed over the entire cluster in a single forward pass. We stack the
cluster's per-sub-chunk query blocks into
$\mathbf{Q}_{\mathcal{C}} = [\mathbf{Q}_{j_1}; \ldots; \mathbf{Q}_{j_{|\mathcal{C}|}}] \in \mathbb{R}^{n_{\mathcal{C}} \times d}$,
where $[\cdot;\cdot]$ is row-wise concatenation and the row count
$n_{\mathcal{C}} = \sum_{j \in \mathcal{C}} |\mathbf{X}_j| \leq |\mathcal{C}| L$
(equality unless $\mathcal{C}$ holds the final, shorter sub-chunk).
The sparse attention output, per
Equation~\eqref{eq:sparse-attention}, is then computed as:
\begin{equation}
    \hat{\mathbf{O}}_{\mathcal{C}}^{(h)} = \text{softmax}\!\left(\frac{\mathbf{Q}_{\mathcal{C}}^{(h)} (\mathbf{K}^{(h)}_{\text{select}})^\top}{\sqrt{d_h}} + \mathbf{M}_{\mathcal{C}}\right) \mathbf{V}^{(h)}_{\text{select}},
    \label{eq:scr-cluster-attention}
\end{equation}
where $\mathbf{K}^{(h)}_{\text{select}}, \mathbf{V}^{(h)}_{\text{select}}$
gather three groups of KV Cache entries: the shared top-$k$ set
$\mathcal{I}_{\mathcal{C}}$ from Equation~\eqref{eq:scr-cluster-selection},
retrieved once from the anchor and reused by every query in the cluster;
the attention
sinks (the first $n_{\text{init}}$ tokens), and a local window. Crucially,
this local window is not fixed at $n_{\text{local}}^{\text{base}}$: it is
expanded to span the entire current cluster, so that every query attends
to all preceding tokens within its own cluster. The anchor-shared
retrieval therefore approximates only the \emph{distant} global context,
whereas all local, intra-cluster attention is computed \emph{exactly}.
After de-duplication, these three groups form
$k_{\text{eff}} \leq k + n_{\text{init}} + n_{\text{local}}$ distinct
entries---where $n_{\text{local}}$ is the cluster-spanning local-window
size---so that
$\mathbf{K}^{(h)}_{\text{select}}, \mathbf{V}^{(h)}_{\text{select}} \in \mathbb{R}^{k_{\text{eff}} \times d_h}$
and $\hat{\mathbf{O}}_{\mathcal{C}}^{(h)} \in \mathbb{R}^{n_{\mathcal{C}} \times d_h}$;
this generalizes Equation~\eqref{eq:sparse-attention}, whose selected set
had exactly $k$ entries. The causal mask
$\mathbf{M}_{\mathcal{C}} \in \mathbb{R}^{n_{\mathcal{C}} \times k_{\text{eff}}}$
sets entry $(r,s)$ to $0$ when selected key $s$ lies at or before the
position of query row $r$, and to $-\infty$ otherwise, so each query
attends only to keys preceding it. The whole cluster is processed in a
single invocation of the FlashInfer paged attention
kernel~\cite{ye2025flashinfer}. Algorithm~\ref{alg:scr-prefill}
summarizes the resulting per-outer-chunk procedure end to end.

\begin{algorithm}[t]
\caption{Similarity Chunk Rolling for one outer chunk.}
\label{alg:scr-prefill}
\begin{algorithmic}[1]
\Require Outer chunk queries $\mathbf{Q} \in \mathbb{R}^{\ell \times d}$; KV Cache $(\mathbf{K}_{\text{cache}}, \mathbf{V}_{\text{cache}})$; threshold $\theta$; sub-chunk size $L$; max cluster size $L_{\max}$; budgets $k_{\text{base}}, n_{\text{local}}^{\text{base}}, n_{\text{init}}$
\State Split $\mathbf{Q}$ into $M = \lceil \ell / L \rceil$ sub-chunks $\{\mathbf{X}_1, \ldots, \mathbf{X}_M\}$
\State Compute normalized fingerprints $\{\tilde{\mathbf{v}}_1, \ldots, \tilde{\mathbf{v}}_M\}$ via Eq.~\eqref{eq:chunk-fingerprint} and L2-normalization
\State $\mathcal{P} \gets \textsc{AnchorClustering}(\{\tilde{\mathbf{v}}_i\}, \theta, L_{\max})$ \Comment{Algorithm~\ref{alg:scr-clustering}}
\ForAll{cluster $\mathcal{C} \in \mathcal{P}$}
    \State $\mathbf{v}_a \gets$ anchor fingerprint of $\mathcal{C}$ (first sub-chunk)
    \State $\mathcal{I}_{\mathcal{C}} \gets \mathcal{S}(\mathbf{v}_a, \mathbf{K}_{\text{cache}}, k_{\text{base}})$ \Comment{Eq.~\eqref{eq:scr-cluster-selection}}
    \State Augment $\mathcal{I}_{\mathcal{C}}$ with attention sinks and a cluster-spanning local window
    \State Compute cluster attention $\hat{\mathbf{O}}_{\mathcal{C}}$ via Eq.~\eqref{eq:scr-cluster-attention} with causal mask
    \State Write $\hat{\mathbf{O}}_{\mathcal{C}}$ to output positions of $\mathcal{C}$ in $\mathbf{O}$
\EndFor
\State \Return $\mathbf{O}$
\end{algorithmic}
\end{algorithm}

\subsection{Self-Improving Agentic Context Engineering}
\label{subsec:agentic-engine}

This subsection presents the second component of SCALE: an agentic 
layer that addresses the \textit{semantic} bottleneck of deploying 
LLMs in region-specific educational settings 
(§\ref{sec:related}, §\ref{subsec:obs-agentic}). Rather than baking 
localized knowledge into model weights through expensive fine-tuning, 
our agentic layer continuously curates an evolving \textit{playbook} 
of curriculum-specific heuristics, common student mistakes, and 
correct teaching trajectories. The playbook is the runtime knowledge 
substrate that grounds each tutoring response, and it is incrementally 
refined through continual interaction with learners.

We structure the agentic layer around three specialized 
roles---Generator, Reflector, and Curator---following the design of 
Agentic Context Engineering (ACE)~\cite{agenticcontext2026}, and 
extend it with domain-specific adaptations for educational tutoring.

\subsubsection{Agentic Architecture Overview}
\label{subsubsec:ace-architecture}

Following the ACE paradigm of Definition~\ref{def:ace}, our agentic layer
keeps the backbone parameters $\phi$ fixed and adapts the playbook
$\mathcal{P}_t$ through the validated update
$\mathcal{P}_{t+1}=\operatorname{Apply}(\mathcal{P}_t,\Delta_t)$
(Equation~\eqref{eq:ace-context-update}; the validated $\operatorname{Apply}$
is made concrete in Equation~\eqref{eq:ace-safe-curation}), where the edit
set $\Delta_t$ is distilled from episode $t$ by the Curator
(§\ref{subsubsec:ace-curator})
and the feedback $y_t$ is a reference answer, an automated grader result,
or a teacher correction. Concretely, a playbook is a set of itemized
entries $b_j=(\mathrm{id}_j,\mathrm{section}_j,\mathrm{text}_j,h_j,a_j)$,
where $h_j$ and $a_j$ count helpful and harmful uses, respectively.

Each episode follows a division of labor inherited from
ACE~\cite{agenticcontext2026}: the Generator $G_\phi$ produces a response
and reasoning trajectory, the Reflector $R_\phi$ diagnoses that trajectory
against the available feedback, and the Curator $C_\phi$ converts the
resulting insight into small, structured edits. Each role is formalized
where it is described, in
§\ref{subsubsec:ace-generator}--§\ref{subsubsec:ace-curator}. All roles
may use the same locally hosted backbone model; their different prompts
and output schemas, rather than access to a stronger teacher model,
create the specialization.

Our contribution extends this base loop with three complementary 
mechanisms: (i) Retrieval-Augmented Execution (RAE), which selects 
relevant rules before generation; (ii) a Failure Memory Bank (FMB), 
which supplies analogous past errors to the Reflector; and (iii) a 
lightweight adversarial curriculum, which periodically creates one 
playbook-conditioned stress test. These mechanisms operate at distinct 
points in the loop and are ablated independently in our experiments.

Algorithm~\ref{alg:scale-agentic} summarizes the complete update loop. 
It supports both offline adaptation over a curated corpus and online 
adaptation after a tutoring interaction. In either setting, the 
playbook and optional failure memory persist across episodes, making 
the learned knowledge easy to audit, edit, and deploy under 
data-sovereignty constraints.

\begin{algorithm}[t]
\caption{Self-improving agentic context engineering.}
\label{alg:scale-agentic}
\begin{algorithmic}[1]
\Require Initial playbook $\mathcal{P}_0$; interaction stream $\mathcal{D}$; optional FMB $\mathcal{F}$; adversarial frequency $f_{\mathrm{adv}}$
\Ensure Updated playbook $\mathcal{P}$ and failure memory $\mathcal{F}$
\For{$(q_t,c_t,y_t) \in \mathcal{D}$}
    \State $\mathcal{P}^{(t)} \gets \operatorname{RetrieveRules}(\mathcal{P},q_t,c_t)$ \Comment{optional RAE}
    \State $(\rho_t,\hat y_t,U_t) \gets G_\phi(q_t,c_t,\mathcal{P}^{(t)})$
    \State $A_t \gets \operatorname{RetrieveFailures}(\mathcal{F},q_t)$ \Comment{optional FMB}
    \State $(r_t,\Gamma_t) \gets R_\phi(q_t,\rho_t,\hat y_t,y_t,U_t,A_t)$
    \State Update helpful/harmful counters of entries in $U_t$ using $\Gamma_t$
    \If{$\hat y_t \neq y_t$ and FMB is enabled}
        \State $\mathcal{F} \gets \mathcal{F} \cup \operatorname{Failure}(q_t,\hat y_t,y_t,r_t)$
    \EndIf
    \State $\Delta_t \gets C_\phi(\mathcal{P},r_t,\Gamma_t)$
    \State $\mathcal{P} \gets \operatorname{Apply}(\mathcal{P},\Delta_t)$
    \If{$t \bmod f_{\mathrm{adv}}=0$}
        \State $\tilde z_t \gets \operatorname{Adversarial}(\mathcal{P},q_t,c_t,y_t)$
        \State Route $\tilde z_t$ through the same Generator--Reflector--Curator loop if it exposes an error
    \EndIf
\EndFor
\State \Return $\mathcal{P},\mathcal{F}$
\end{algorithmic}
\end{algorithm}

\subsubsection{Generator - Producing Tutoring Trajectories}
\label{subsubsec:ace-generator}

The Generator is prompted with the student query, instructional
context, and a structured subset of the current playbook. Formally, it
returns a tutoring trajectory $\rho_t$, an answer $\hat y_t$, and the
identifiers $U_t$ of the entries used in its reasoning,
\begin{equation}
    (\rho_t,\hat{y}_t,U_t) = G_\phi(q_t,c_t,\mathcal{P}_t),
    \label{eq:ace-generator-role}
\end{equation}
where $\mathcal{P}_t$ is the playbook supplied to the call---the focused
subset selected by RAE (Equation~\eqref{eq:ace-rae},
§\ref{subsubsec:ace-update}) when it is enabled, otherwise the full
playbook. Requiring entry
identifiers makes playbook use observable: the subsequent Reflector can
attribute a successful or failed response to concrete advice rather
than treating the context as an opaque prompt.

For tutoring, $\rho_t$ can include a worked explanation, a Socratic 
hint sequence, or a tool-use trace, depending on the learning objective. 
The response contract separates pedagogical form from factual content: 
the Generator must answer the question while exposing enough intermediate 
reasoning for feedback to identify whether the failure arose from a 
misunderstood local fact, an inappropriate instructional strategy, or a 
missing prerequisite. This trajectory is not retained as hidden model 
state; it is transient evidence for reflection and curation.

\subsubsection{Reflector - Extracting Insights from Execution Feedback}
\label{subsubsec:ace-reflector}

The Reflector receives the student query $q_t$; the trajectory $\rho_t$,
answer $\hat y_t$, and cited entries $U_t$ produced by the Generator
(Equation~\eqref{eq:ace-generator-role}); and the strongest feedback
signal $y_t$ available. During
offline training this signal can be a reference answer or deterministic
grader; during deployment it can be a teacher correction, tool result,
or explicit student follow-up. Formally, it maps these inputs to a
structured reflection $r_t$ and a set of tags $\Gamma_t$ over the cited
entries,
\begin{equation}
    (r_t,\Gamma_t) = R_\phi(q_t,\rho_t,\hat{y}_t,y_t,U_t),
    \label{eq:ace-reflector-role}
\end{equation}
where each cited entry receives a tag in
$\{\textsc{helpful},\textsc{neutral},\textsc{harmful}\}$ and the
reflection itself is
\begin{equation}
    r_t=(e_t,z_t,a_t^{\star},\kappa_t),
    \label{eq:ace-reflection}
\end{equation}
where $e_t$ identifies the observed error, $z_t$ explains its root
cause, $a_t^{\star}$ specifies a correct approach, and $\kappa_t$ states a
reusable key insight.

This explicit diagnosis is important for localized tutoring. For example, 
an answer can be linguistically fluent yet fail because it uses an 
incorrect regional historical fact, conflates two curriculum definitions, 
or gives a pedagogically unsuitable shortcut. The Reflector converts such 
instance-level feedback into an actionable explanation that can be reused 
without exposing the original student's private record. Multiple 
reflection rounds may be used when a first diagnosis is incomplete; the 
final reflection is passed to the Curator. Because a reflection is later 
converted into persistent context, the Reflector is also the critical 
point at which hallucinated diagnosis can enter the long-horizon learning 
loop. We therefore preserve the evidence, cited bullet identifiers, and 
the resulting delta as auditable artifacts, rather than treating a 
reflection as an untraceable free-form summary.

\paragraph{Failure Memory Bank}
To support analogical reflection, SCALE maintains a Failure Memory Bank
(FMB), separate from the rule-oriented playbook. A memory record pairs the
failed question, prediction, and reference outcome with the error $e_i$,
root cause $z_i$, and key insight $\kappa_i$ distilled by the reflection
(Equation~\eqref{eq:ace-reflection}):
\begin{equation}
    f_i=(q_i,\hat y_i,y_i,e_i,z_i,\kappa_i,\mathbf{e}_i),
    \label{eq:ace-fmb-schema}
\end{equation}
where $\mathbf{e}_i=\operatorname{norm}(E(q_i))$ is an in-memory
embedding used only for retrieval. A record is added after reflection 
when the observed answer disagrees with the available reference outcome,
so the record contains distilled diagnosis rather than an unprocessed 
trajectory. For a new query $q$, normalized embedding vectors retrieve 
the $K$ closest prior failures:
\begin{equation}
    \mathcal{N}_K(q)=\operatorname{TopK}_{f_i\in\mathcal{F}}
    \frac{E(q)^\top E(q_i)}{\|E(q)\|_2\|E(q_i)\|_2}.
    \label{eq:ace-fmb-retrieval}
\end{equation}
The retrieved cases are rendered as additional context for the Reflector, 
allowing it to recognize recurring failure patterns. FMB therefore does 
not replace the playbook: the playbook guides execution, whereas FMB 
supplies episodic evidence to improve diagnosis.

\subsubsection{Curator - Incremental Playbook Updates}
\label{subsubsec:ace-curator}

The Curator turns the reflection and entry tags $(r_t,\Gamma_t)$ of
Equation~\eqref{eq:ace-reflector-role} into a compact set of typed delta
operations $\Delta_t$, rather than rewriting the full prompt:
\begin{equation}
    \Delta_t = C_\phi(\mathcal{P}_t,r_t,\Gamma_t).
    \label{eq:ace-curator-role}
\end{equation}
The operational path
primarily uses an \textsc{Add} request over a named playbook section;
the surrounding schema also accommodates targeted update, merge, and 
delete operations as the playbook maintenance policy evolves. The JSON 
operation schema is validated before application; malformed output yields 
an empty delta and leaves the existing playbook unchanged. This safety
guard is the concrete definition of the $\operatorname{Apply}$ operator
of Equation~\eqref{eq:ace-context-update}:
\begin{equation}
    \operatorname{Apply}(\mathcal{P}_t,\Delta_t)=\begin{cases}
    \Phi(\mathcal{P}_t,\Delta_t), & \text{if }\Delta_t\text{ is valid},\\
    \mathcal{P}_t, & \text{otherwise},
    \end{cases}
    \label{eq:ace-safe-curation}
\end{equation}
where $\Phi$ is deterministic, non-LLM application logic. Periodic
merging of redundant entries (the ``refine'' step of the grow-and-refine
policy) is applied separately and is described below.

Entries are organized by reusable strategies, common mistakes, worked 
templates, and troubleshooting guidance. The Curator follows a 
grow-and-refine policy: it appends a new entry when the reflection exposes 
a genuinely absent rule; otherwise it updates usage counters or merges 
overlapping entries. This localized update mechanism preserves prior 
knowledge and prevents the context-collapse failure mode associated with 
repeated monolithic prompt rewrites.

\subsubsection{Incremental Updates and Context Collapse Prevention}
\label{subsubsec:ace-update}

As the playbook grows, passing every entry to every generation is both 
inefficient and vulnerable to distraction from irrelevant rules. We use 
Retrieval-Augmented Execution (RAE) to select only the most relevant 
playbook bullets before invoking the Generator. Each entry
$b_j=(\mathrm{id}_j,\mathrm{section}_j,\mathrm{text}_j,h_j,a_j)$ from the
schema defined above contributes its bullet text $\mathrm{text}_j$,
which is embedded and L2-normalized.
The implementation parses the structured playbook, indexes all bullet
texts with FAISS inner-product search, and retains section headers when
reconstructing a focused prompt. For query context $u_t=q_t\oplus c_t$,
RAE computes
\begin{align}
    \mathbf{e}_j &= \operatorname{norm}(E(\mathrm{text}_j)),\qquad
    \mathbf{e}_q=\operatorname{norm}(E(u_t)),\\
    I_K &= \operatorname{TopK}_j\; \mathbf{e}_q^\top\mathbf{e}_j,
    \label{eq:ace-rae}
\end{align}
and supplies the focused playbook
$\mathcal{P}^{(t)}=\{b_j\mid j\in I_K\}$ to the Generator. Section 
headers are retained so that retrieved bullets preserve their original 
organizational context. The index is rebuilt after each accepted Curator 
update; if the index is unavailable or the playbook contains at most $K$ 
entries, the system safely falls back to the full playbook.

RAE is deliberately placed on the execution path, rather than used only 
for post-hoc analysis: the same retrieval call conditions initial 
generation, post-reflection regeneration, post-Curator generation, 
validation, and adversarial execution. It therefore reduces both prompt 
noise and inference cost whenever the playbook becomes long, while entry 
identifiers still allow the Reflector to assign credit or blame to the 
rules actually used.

RAE and FMB share the same embedding encoder when both are enabled, 
avoiding a duplicate model copy on memory-constrained hardware. RAE 
retrieves rules for the Generator, while FMB retrieves past failures for 
the Reflector; this separation prevents a past erroneous trajectory from 
being directly treated as an execution instruction.

To bound long-run context growth, SCALE periodically compares entry 
embeddings and merges semantically redundant entries when their cosine 
similarity exceeds a configurable threshold. Pruning or deduplication is 
performed only through validated, localized operations. Consequently, an 
update to one educational rule neither discards unrelated rules nor 
requires re-generating the entire playbook.

\subsubsection{Adversarial Curriculum}
\label{subsubsec:ace-adversarial}

We use a lightweight adversarial curriculum during adaptation to
expose playbook weaknesses not covered by routine interactions. For the
Formula and FiNER experiments it is invoked when
$t\bmod f_{\mathrm{adv}}=0$; for AppWorld it runs after the original
task trajectory has completed. A single adversarial 
LLM call, conditioned on the current playbook and a recent training 
sample, returns exactly one plausible challenge
\begin{equation}
    \tilde z_t=(\tilde q_t,\tilde c_t,\tilde y_t,\tilde a_t,\tilde v_t),
    \label{eq:ace-adversarial-schema}
\end{equation}
containing a question, minimal context, target answer, attack rationale, 
and vulnerability hint. The standard Generator executes this challenge 
under the current playbook and the data processor compares its answer 
against the generated target. If $\hat{\tilde y}_t\neq\tilde y_t$, the
resulting trajectory is routed through the same
Generator--Reflector--Curator loop as a genuine failure
(Equations~\eqref{eq:ace-generator-role}--\eqref{eq:ace-curator-role}), and
can also enter the FMB after reflection, letting a synthetically exposed
weakness inform later analogical diagnosis. The
adversarial component therefore creates additional training signals, but 
never edits the playbook directly. Since this mechanism uses one
generated target without an independent verifier, we treat it as active 
stress testing rather than formal evidence of an error. This distinction 
is essential: the contribution is a low-cost curriculum generator, not a
claim of adversarial verification.

\section{Experiments and Results}
\label{sec:experiments}

We evaluate SCALE in three parts. Because the two components target
different bottlenecks and call for different baselines, we first study
each in isolation: the long-context inference engine
(§\ref{subsec:exp-scr}) is measured against a state-of-the-art selective
sparse attention method on long-context benchmarks, whereas the
self-improving agentic layer (§\ref{subsec:exp-agentic}) is evaluated as
a controlled ablation over its mechanisms on agentic reasoning
benchmarks. We then measure the two together in the configuration DeepEdu
actually deploys---the agentic layer running on top of SCR
(§\ref{subsec:exp-scr-agent}).

\subsection{Long-Context Inference Engine}
\label{subsec:exp-scr}

We evaluate the long-context inference engine of SCALE
(§\ref{subsec:long-context-engine}) against a strong selective sparse
attention baseline, TokenSelect~\cite{tokenselect2025}, on two 
established long-context benchmarks. We aim to answer three research 
questions:
\begin{itemize}
    \item \textbf{RQ1}: Does Similarity Chunk Rolling (SCR) reduce 
    prefill latency compared to per-sub-chunk selective sparse 
    attention?
    \item \textbf{RQ2}: Does the cluster-level retrieval paradigm 
    preserve or improve task accuracy compared to per-sub-chunk 
    retrieval?
    \item \textbf{RQ3}: How sensitive is SCR to its hyperparameters, 
    specifically the similarity threshold $\theta$ and the maximum 
    cluster size $L_{\max}$?
\end{itemize}

\subsubsection{Experimental Setup}
\label{subsec:exp-setup}

\paragraph{Hardware and software} 
All experiments are conducted on a single NVIDIA H100 GPU with
80\,GB HBM3. We implement SCR on the SGLang~\cite{zheng2024sglang}
serving framework with the FlashInfer~\cite{ye2025flashinfer} attention
backend, in full bfloat16 precision. The backbone model 
is Qwen2-7B-Instruct~\cite{yang2024qwen2}, configured with a maximum 
context length of $1{,}048{,}576$ tokens, RoPE base $10^6$, outer 
chunk size $L_{\text{outer}} = 8192$, and inner sub-chunk size 
$L = 512$.

\paragraph{Baseline} 
Our primary baseline is \textbf{TokenSelect}~\cite{tokenselect2025}, 
the current state-of-the-art query-aware token-level selective sparse 
attention method. TokenSelect invokes its selection function 
$\mathcal{S}$ independently per sub-chunk during prefill, providing 
a direct comparison point against SCR's cluster-level amortization. 
We use the same selection function (head soft-voting), retrieval
budget ($k_{\text{base}} = 8192$), base local window
($n_{\text{local}}^{\text{base}} = 512$), and attention sinks
($n_{\text{init}} = 128$) for both methods. By construction, however,
SCR's local window expands to span each cluster
(§\ref{subsubsec:scr-retrieval}), so the comparison reflects the full
cluster-level mechanism: amortized retrieval, which drives the latency
savings, together with the wider cluster-spanning local attention,
which drives the accuracy gains on structured-retrieval tasks.

\paragraph{SCR configurations}
We evaluate SCR across the full grid of two hyperparameters:
similarity threshold $\theta \in \{0.95, 0.97, 0.99\}$ and maximum
cluster size $L_{\max} \in \{1024, 2048, 4096\}$ (i.e., $2$, $4$, and
$8$ sub-chunks per cluster), yielding nine configurations.

\paragraph{Benchmarks} 
We use two standard long-context benchmarks:
\begin{itemize}
    \item \textbf{InfiniteBench}~\cite{zhang_etal_2024_bench}: six tasks
    spanning code debugging (Code.D), structured retrieval
    (key--value retrieval R.KV, number string R.Num, passkey R.PK),
    narrative dialogue QA (En.Dia), and math find (Math.F); we use these
    abbreviations, following~\cite{tokenselect2025}, throughout. Each
    task is evaluated on its full test set with native long-context
    inputs.
    \item \textbf{RULER}~\cite{ruler2024}: eight needle-in-a-haystack
    (NIAH) tasks---single-, multi-key, multi-value, and multi-query
    variants---at $100$ samples per task across six context lengths from
    $4096$ to $131{,}072$ tokens, isolating retrieval accuracy as
    context length grows.
\end{itemize}
We report average accuracy across tasks and average time-to-first-token 
(TTFT) latency in seconds, with TTFT measured from the start of 
prefill to the emission of the first decoded token.


\subsubsection{Accuracy and Latency on InfiniteBench}
\label{subsec:exp-infinitebench}

\paragraph{Accuracy}
Table~\ref{tab:infbench-accuracy} reports per-task accuracy on
representative InfiniteBench tasks (the full six-task grid is in
Appendix~\ref{app:infbench-full}). SCR matches or exceeds the
TokenSelect baseline on nearly all tasks. Most notably, on R.KV---a task
with long structured contexts that historically stresses sparse
attention---SCR with $L_{\max} = 4096$ reaches $98.0\%$ versus the
baseline's $91.0\%$, a \textbf{$7$-point absolute improvement}, and
this accuracy grows monotonically with cluster size
($94.0 \to 96.2 \to 98.0\%$ as $L_{\max}$ increases from $1024$ to
$4096$). We attribute this gain not to amortization---which affects only
latency---but to the wider effective local window that cluster-level
processing yields as an intrinsic byproduct, preserving more local
context relevant to structured retrieval. Crucially, SCR obtains this
wider window \emph{while} amortizing retrieval; a baseline could match
it only by widening its own per-sub-chunk attention, at the cost of the
latency advantage. SCR also edges out the
baseline on Code.D ($28.8\%$ vs.\ $26.7\%$) and
En.Dia ($11.5\%$ vs.\ $10.0\%$), while
R.Num and R.PK are saturated at
$100\%$ for both TokenSelect and every SCR configuration. Against the
other long-context baselines the margin on structured retrieval is far
larger: the strongest competing sparse method on R.KV
(NTK) reaches only $59.8\%$, and memory-based methods such as InfLLM
fall below $6\%$, underscoring how much cluster-level selection
preserves.

\begin{table}[t]
\centering
\caption{Per-task accuracy (\%) on representative InfiniteBench tasks
(Qwen2-7B-Instruct). R.Num and R.PK are
saturated at $100\%$ and omitted here; the full six-task,
nine-configuration grid is in Table~\ref{tab:infbench-accuracy-full}.
Best per column in bold.}
\label{tab:infbench-accuracy}
\small
\setlength{\tabcolsep}{5pt}
\begin{tabular}{lcccc}
\toprule
\textbf{Method} & \textbf{Code.D} & \textbf{R.KV} & \textbf{En.Dia} & \textbf{Math.F} \\
\midrule
Qwen2-7B-Instruct~\cite{yang2024qwen2}          & 28.2 & 19.0 & \phantom{0}8.5 & 19.7 \\
NTK~\cite{ntk2023web}                  & 24.9 & 59.8 & \phantom{0}7.5 & \textbf{27.7} \\
SelfExtend~\cite{jin2024selfextend}    & \phantom{0}8.1 & \phantom{0}0.0 & \phantom{0}5.0 & \phantom{0}2.3 \\
StreamingLLM~\cite{xiao2024efficient}  & 27.9 & \phantom{0}2.4 & \phantom{0}3.5 & 19.4 \\
InfLLM~\cite{xiao2024infllm}           & 27.4 & \phantom{0}5.4 & \phantom{0}7.5 & 24.0 \\
TokenSelect~\cite{tokenselect2025}     & 26.7 & 91.0 & 10.0 & 24.0 \\
\midrule
$\theta\!=\!0.95$, $L_{\max}\!=\!1024$      & 27.6 & 94.0 & 10.0 & 23.7 \\
$\theta\!=\!0.95$, $L_{\max}\!=\!2048$      & 26.1 & 96.2 & 10.0 & 23.4 \\
$\theta\!=\!0.95$, $L_{\max}\!=\!4096$      & 24.9 & \textbf{98.0} & \phantom{0}6.0 & 22.6 \\
$\theta\!=\!0.99$, $L_{\max}\!=\!1024$      & \textbf{28.8} & 94.0 & \textbf{11.5} & 23.1 \\
\bottomrule
\end{tabular}
\end{table}

\paragraph{Latency}
Table~\ref{tab:infbench-latency} reports TTFT latency per task. SCR
delivers substantial latency reductions over the baseline on all
tasks. Averaging across tasks, SCR with $\theta = 0.95$ and
$L_{\max} = 4096$ achieves roughly a \textbf{$35\%$ latency
reduction}: from $13.2$\,s to $8.3$\,s on R.KV,
from $8.8$\,s to $5.7$\,s on Code.D, and from
$8.7$\,s to $5.7$\,s on Math.F.

\begin{table}[t]
\centering
\caption{Per-task TTFT latency (s) on representative InfiniteBench tasks.
SCR delivers roughly a $35\%$ latency reduction over the TokenSelect
baseline (top-$k = 8192$), with the largest gains on long-context
retrieval tasks. The complete nine-configuration grid over all six tasks
appears in Table~\ref{tab:infbench-latency-full}. Best (lowest) per
column in bold.}
\label{tab:infbench-latency}
\small
\setlength{\tabcolsep}{5pt}
\begin{tabular}{lcccc}
\toprule
\textbf{Method} & \textbf{Code.D} & \textbf{R.KV} & \textbf{En.Dia} & \textbf{Math.F} \\
\midrule
TokenSelect                            & 8.8 & 13.2 & 7.7 & 8.7 \\
\midrule
$\theta\!=\!0.95$, $L_{\max}\!=\!1024$      & 6.7 & \phantom{0}9.9 & 6.0 & 6.6 \\
$\theta\!=\!0.95$, $L_{\max}\!=\!2048$      & 5.9 & \phantom{0}8.7 & 5.3 & 5.9 \\
$\theta\!=\!0.95$, $L_{\max}\!=\!4096$      & \textbf{5.7} & \phantom{0}\textbf{8.3} & \textbf{5.1} & \textbf{5.7} \\
$\theta\!=\!0.99$, $L_{\max}\!=\!1024$      & 7.3 & \phantom{0}9.9 & 6.2 & 6.7 \\
\bottomrule
\end{tabular}
\end{table}


\paragraph{Where the speedup comes from}
To see \emph{why} SCR is faster, we instrument the prefill pass and
attribute wall-clock time to each stage
(Figure~\ref{fig:prefill-breakdown}, R.KV, $n\!=\!30$). Almost the entire
gain comes from a single stage---token retrieval---which falls from $4.5$
to $0.9$\,s per sample ($77\%$ of the total saving), because SCR issues
$7.7\times$ fewer retrieval calls ($8{,}959\!\to\!1{,}165$), close to the
analytic $8\times$ bound (§\ref{subsubsec:scr-clustering}). Attention
compute barely changes ($-4\%$) even though SCR processes $5.5\times$
fewer KV tokens: the win is in how \emph{often} the model selects tokens,
not in the attention FLOPs. As retrieval cost grows with context length,
this saving only widens for the longer contexts that motivate SCR.

\begin{figure}[t]
\centering
\includegraphics[width=0.96\columnwidth]{figures/prefill_breakdown.png}
\caption{Per-sample prefill time by stage on InfiniteBench R.KV
(Qwen2-7B-Instruct; separate instrumented run, $n\!=\!30$). Overall TTFT falls $13.1$\,s $\to$
$8.3$\,s ($-37\%$); almost the entire reduction comes from the
token-retrieval stage, while attention compute is nearly unchanged.}
\label{fig:prefill-breakdown}
\end{figure}


\subsubsection{Scaling with Context Length on RULER}
\label{subsec:exp-ruler}

\paragraph{Accuracy} 
Table~\ref{tab:ruler-accuracy} reports average accuracy across the 
eight RULER tasks at six context lengths. SCR variants match or 
slightly exceed the TokenSelect baseline at every context length, 
confirming that cluster-level amortization does not compromise 
retrieval fidelity. At the longest context of $128$K tokens, the
best SCR configuration ($L_{\max} = 4096$) achieves
$65.28\%$ accuracy versus the baseline's $64.78\%$, a small but
consistent improvement.

\begin{table}[t]
\centering
\caption{Average accuracy (\%) on RULER across eight 
needle-in-a-haystack tasks at varying context lengths. SCR variants 
preserve or slightly improve over the TokenSelect baseline at every 
context length.}
\label{tab:ruler-accuracy}
\footnotesize
\setlength{\tabcolsep}{4pt}
\begin{tabular}{lcccccc}
\toprule
\textbf{Method} & \textbf{4K} & \textbf{8K} & \textbf{16K} & \textbf{32K} & \textbf{64K} & \textbf{128K} \\
\midrule
TokenSelect                            & 96.53 & 97.72 & 95.16 & \textbf{78.69} & 69.09 & 64.78 \\
\midrule
$\theta\!=\!0.95$, $L_{\max}\!=\!1024$      & \textbf{96.75} & 97.59 & 95.28 & 78.53 & 68.66 & 64.88 \\
$\theta\!=\!0.95$, $L_{\max}\!=\!4096$      & 96.69 & \textbf{97.81} & \textbf{95.97} & 77.81 & \textbf{69.72} & \textbf{65.28} \\
\bottomrule
\end{tabular}
\end{table}

\paragraph{Latency} 
Table~\ref{tab:ruler-latency} reports average TTFT in seconds. SCR 
delivers progressively larger latency savings as context length grows. 
At $4$K--$8$K tokens, SCR variants and baseline are within $5\%$ of 
each other (both near $0.5$--$0.7$\,s), as retrieval cost is 
negligible at short contexts. At $128$K tokens, SCR with 
$L_{\max} = 4096$ achieves $6.64$\,s versus the baseline's $9.89$\,s, 
a \textbf{$32.9\%$ latency reduction}. The advantage widens
super-linearly with context length---$3.0\%$ at $16$K, $16.7\%$ at
$32$K, $25.9\%$ at $64$K, and $32.9\%$ at $128$K---because longer
prefills process more outer chunks, creating more clustering
opportunities and a growing cumulative reduction in retrieval
invocations. This confirms that the amortization predicted by
§\ref{sec:motivations} materializes in wall-clock latency on real
hardware.

\begin{table}[t]
\centering
\caption{Average TTFT latency (s) on RULER at varying context lengths. 
SCR achieves substantial latency reduction at long contexts, with the 
gap widening as context length grows.}
\label{tab:ruler-latency}
\footnotesize
\setlength{\tabcolsep}{4pt}
\begin{tabular}{lcccccc}
\toprule
\textbf{Method} & \textbf{4K} & \textbf{8K} & \textbf{16K} & \textbf{32K} & \textbf{64K} & \textbf{128K} \\
\midrule
TokenSelect                            & \textbf{0.52} & 0.65 & 1.94 & 3.17 & 5.27 & 9.89 \\
\midrule
$\theta\!=\!0.95$, $L_{\max}\!=\!1024$      & 0.53 & 0.65 & 1.88 & 2.94 & 4.57 & 8.16 \\
$\theta\!=\!0.95$, $L_{\max}\!=\!4096$      & \textbf{0.52} & \textbf{0.62} & \textbf{1.75} & \textbf{2.64} & \textbf{3.90} & \textbf{6.64} \\
\bottomrule
\end{tabular}
\end{table}


\subsubsection{Ablation: Effect of $L_{\max}$ and Threshold $\theta$}
\label{subsec:exp-ablation}

\paragraph{Effect of maximum cluster size $L_{\max}$}
Increasing $L_{\max}$ from $1024$ to $4096$ at fixed $\theta = 0.95$
monotonically reduces latency---on R.KV, TTFT drops
from $9.9$\,s to $8.3$\,s---confirming that larger clusters yield more
amortization. The accuracy effect is task-dependent: structured
retrieval benefits from larger clusters (R.KV
climbs from $94.0\%$ to $98.0\%$ as $L_{\max}$ grows from $1024$ to
$4096$), whereas code and dialogue reasoning slightly favor smaller
clusters (Code.D and En.Dia
peak at $L_{\max} = 1024$). This asymmetry reflects the trade-off
between amortization gain and selection granularity: when clusters
become large, the single anchor fingerprint $\mathbf{v}_a$ must stand
in for an increasingly wide span of sub-chunks, so members further
from the anchor may carry retrieval needs it does not fully capture.
This helps tasks with locally coherent context but mildly hurts tasks
requiring fine-grained per-sub-chunk selection. We adopt $L_{\max} = 4096$ as default for its consistent
latency advantage and strong structured-retrieval accuracy, noting
that practitioners may reduce $L_{\max}$ when fine-grained reasoning
is paramount.

\paragraph{Effect of similarity threshold $\theta$}
Sweeping $\theta \in \{0.95, 0.97, 0.99\}$ at fixed $L_{\max}$
(full grid in Appendix~\ref{app:infbench-full}) reveals
only minor accuracy variation ($\pm 2$ points) with near-identical
latency, consistent with the cluster size distribution observed in
§\ref{subsec:obs-cluster}: across this range the majority of
clustering decisions still reach the maximum cluster size, so the
effective amortization is largely insensitive to $\theta$. We adopt
$\theta = 0.95$ as default for its slightly more aggressive
amortization, with higher thresholds available when finer cluster
granularity is desired.

\subsection{Self-Improving Agentic Layer}
\label{subsec:exp-agentic}

We evaluate the self-improving agentic layer as a controlled ablation
over its three added mechanisms---Retrieval-Augmented Execution (RAE),
the Failure Memory Bank (FMB), and the adversarial curriculum---%
independently of the TokenSelect comparison. We aim to answer three
research questions:
\begin{itemize}
    \item \textbf{RQ4}: Do RAE, FMB, and the adversarial
    curriculum improve grounded financial reasoning on Formula and
    FiNER under the same offline-learning budget?
    \item \textbf{RQ5}: Do these mechanisms transfer to interactive
    AppWorld agents, including both normal and challenge trajectories?
    \item \textbf{RQ6}: As selective retrieval, failure memory, and
    self-initiated stress testing are added incrementally, how does each
    mechanism contribute, and where does it introduce task-specific
    trade-offs?
\end{itemize}

\subsubsection{Experimental Setup}
\label{subsec:agentic-experiments}

\paragraph{Setup and baseline}
All Generator, Reflector, and Curator roles use Qwen3-4B-Instruct-2507,
served through vLLM on the same single H100 80\,GB GPU as the SCR
experiments. Supervision comes only from Ground Truth labels or
deterministic environment outcomes---no configuration learns from
unverified live student interactions---reflecting the intended
educational workflow, in which the system first curates reliable
knowledge from verified data and then deploys the resulting Playbook. Our
baseline is the original ACE architecture~\cite{agenticcontext2026}, a
Generator--Reflector--Curator loop that accumulates a structured Playbook
from Ground Truth feedback; the \textbf{Origin} row in
Tables~\ref{tab:agentic-finance-ablation} and
\ref{tab:agentic-appworld-ablation} is this baseline with RAE, FMB, and
the adversarial curriculum disabled. Every variant shares Origin's
backbone, supervision, seed, token budgets, and evaluation protocol, so
the ablation isolates the effect of the three added mechanisms.

\paragraph{Training configuration for Formula and FiNER}
We train for one epoch with at most three agentic
rounds per instance, a $4{,}000$-token Playbook budget, and a
$2{,}048$-token generation limit. We use seed $42$, evaluate every
$50$ steps, save every $25$ steps, and run evaluation with five workers.
RAE and FMB each retrieve the top $K=10$ entries. The
adversarial curriculum is scheduled every $f_{\mathrm{adv}}=10$ update
steps. Thus, the ablation changes only the switches
$(\mathrm{RAE},\mathrm{FMB},\mathrm{Adv})$; all backbone, optimization,
and evaluation settings are held fixed.

\paragraph{AppWorld protocol and reported metrics}
Among the ACE benchmarks, AppWorld is the closest analogue to our target
setting of AI-assisted learning: it requires an agent to carry out
interactive, multi-step, tool-augmented tasks in a stateful environment,
recovering from its own execution errors along the way---exactly the
mode in which a tutoring agent must operate on a learner's behalf, as
opposed to answering a single self-contained question. Its normal and
challenge splits therefore serve as the shared testbed for both SCALE
components in this setting: they let us probe whether the self-improving
loop improves task success (§\ref{subsec:exp-agentic}) and whether the
SCR engine sustains low latency (§\ref{subsec:exp-scr-agent}), precisely
where educational deployment is hardest---sustained, error-prone,
multi-turn interaction.
For AppWorld we use the same Qwen3-4B-Instruct-2507 backbone and seed
$42$, with RAE and FMB again configured at $K=10$. The environment's
task completion signal is treated as Ground Truth. Unlike the
step-scheduled Formula/FiNER protocol, an AppWorld adversarial case is
generated only after the original task trajectory finishes. This ensures
that the synthetic case probes a completed, observable trajectory and
does not replace or contaminate the original task evaluation. We report
final test accuracy on Formula and FiNER, and Task Goal Completion (TGC)
and Scenario Goal Completion (SGC) on the AppWorld normal and challenge
splits.

% ---- Old full 2^3 combinatorial table (kept for reference) ----
\iffalse
\begin{table*}[t]
\centering
\caption{Offline adaptation on Formula and FiNER using Ground Truth
labels. All runs use Qwen3-4B-Instruct-2507, seed 42, and identical
training settings. Average is the mean of the two accuracies.}
\small
\setlength{\tabcolsep}{8pt}
\begin{tabular}{l|ccc|ccc}
\toprule
Method & RAE & FMB & Adv. & FiNER Acc$\uparrow$ & Formula Acc$\uparrow$ & Average$\uparrow$ \\
\midrule
\rowcolor{gray!18}\multicolumn{7}{c}{Qwen3-4B-Instruct-2507 as Base LLM} \\
Origin & -- & -- & -- & 51.3 & 70.0 & 60.7 \\
\midrule
\rowcolor{gray!18}\multicolumn{7}{c}{Offline Adaptation (Ground Truth)} \\
RAE & \checkmark & -- & -- & 47.6 & 77.5 & 62.6 \\
FMB & -- & \checkmark & -- & \textbf{60.4} & 74.0 & \textbf{67.2} \\
RAE + FMB & \checkmark & \checkmark & -- & 50.0 & 77.0 & 63.5 \\
Adversarial & -- & -- & \checkmark & 47.6 & 74.0 & 60.8 \\
Adversarial + FMB & -- & \checkmark & \checkmark & 55.8 & 66.5 & 61.2 \\
RAE + Adversarial & \checkmark & -- & \checkmark & 52.1 & \textbf{79.5} & 65.8 \\
RAE + FMB + Adversarial & \checkmark & \checkmark & \checkmark & 52.9 & 73.0 & 63.0 \\
\bottomrule
\end{tabular}
\end{table*}
\fi

\begin{table*}[t]
\centering
\caption{Incremental ablation of the agentic mechanisms on Formula and
FiNER using Ground Truth labels. Mechanisms are added cumulatively on top
of the Origin Playbook; all runs use Qwen3-4B-Instruct-2507, seed 42, and
identical training settings. Average is the mean of the two accuracies;
the final row is the full configuration DeepEdu deploys. Best per column in
bold.}
\label{tab:agentic-finance-ablation}
\small
\setlength{\tabcolsep}{8pt}
\begin{tabular*}{0.7\textwidth}{@{\extracolsep{\fill}}l|ccc|ccc@{}}
\toprule
Method & RAE & FMB & Adv. & FiNER Acc$\uparrow$ & Formula Acc$\uparrow$ & Average$\uparrow$ \\
\midrule
Origin & -- & -- & -- & 51.3 & 70.0 & 60.7 \\
\midrule
\quad + RAE & \checkmark & -- & -- & 47.6 & 77.5 & 62.6 \\
\quad + Adversarial & \checkmark & -- & \checkmark & 52.1 & \textbf{79.5} & \textbf{65.8} \\
\quad + FMB (DeepEdu) & \checkmark & \checkmark & \checkmark & \textbf{52.9} & 73.0 & 63.0 \\
\bottomrule
\end{tabular*}
\end{table*}

% ---- Old full 2^3 combinatorial table (kept for reference) ----
\iffalse
\begin{table*}[t]
\centering
\caption{Offline adaptation on AppWorld using Ground Truth environment
outcomes. TGC and SGC are reported for the normal and challenge splits;
Average is the mean of the four reported AppWorld metrics.}
\small
\setlength{\tabcolsep}{6pt}
\begin{tabular}{l|ccc|cc|cc|c}
\toprule
Method & RAE & FMB & Adv. & \multicolumn{2}{c|}{Test-Normal} & \multicolumn{2}{c|}{Test-Challenge} & Average \\
 &  &  &  & TGC$\uparrow$ & SGC$\uparrow$ & TGC$\uparrow$ & SGC$\uparrow$ & $\uparrow$ \\
\midrule
\rowcolor{gray!18}\multicolumn{9}{c}{Qwen3-4B-Instruct-2507 as Base LLM} \\
Origin & -- & -- & -- & 20.8 & 8.9 & 6.7 & 0.7 & 9.3 \\
\midrule
\rowcolor{gray!18}\multicolumn{9}{c}{Offline Adaptation (Ground Truth)} \\
RAE & \checkmark & -- & -- & 20.2 & \textbf{12.5} & 8.2 & 0.7 & 10.4 \\
FMB & -- & \checkmark & -- & 23.2 & \textbf{12.5} & \textbf{8.9} & \textbf{2.9} & \textbf{11.9} \\
RAE + FMB & \checkmark & \checkmark & -- & 21.4 & 10.7 & 7.4 & 2.2 & 10.4 \\
Adversarial & -- & -- & \checkmark & 22.6 & 7.1 & 6.2 & 2.2 & 9.5 \\
Adversarial + FMB & -- & \checkmark & \checkmark & \textbf{24.4} & \textbf{12.5} & 8.6 & 0.7 & 11.6 \\
RAE + Adversarial & \checkmark & -- & \checkmark & 20.2 & 7.1 & 6.2 & 2.2 & 8.9 \\
RAE + FMB + Adversarial & \checkmark & \checkmark & \checkmark & 22.0 & 10.7 & 7.4 & 1.4 & 10.4 \\
\bottomrule
\end{tabular}
\end{table*}
\fi

\begin{table*}[t]
\centering
\caption{Incremental ablation of the agentic mechanisms on AppWorld using
Ground Truth environment outcomes. Mechanisms are added cumulatively on
top of Origin; TGC and SGC are reported for the normal and challenge
splits, and Average is the mean of the four reported metrics. The final
row is the full configuration DeepEdu deploys. Best per column in bold.}
\label{tab:agentic-appworld-ablation}
\small
\setlength{\tabcolsep}{6pt}
\begin{tabular*}{0.85\textwidth}{@{\extracolsep{\fill}}l|ccc|cc|cc|c@{}}
\toprule
Method & RAE & FMB & Adv. & \multicolumn{2}{c|}{Test-Normal} & \multicolumn{2}{c|}{Test-Challenge} & Average \\
 &  &  &  & TGC$\uparrow$ & SGC$\uparrow$ & TGC$\uparrow$ & SGC$\uparrow$ & $\uparrow$ \\
\midrule
Origin & -- & -- & -- & 20.8 & 8.9 & 6.7 & 0.7 & 9.3 \\
\midrule
\quad + RAE & \checkmark & -- & -- & 20.2 & \textbf{12.5} & \textbf{8.2} & 0.7 & \textbf{10.4} \\
\quad + Adversarial & \checkmark & -- & \checkmark & 20.2 & 7.1 & 6.2 & \textbf{2.2} & 8.9 \\
\quad + FMB (DeepEdu) & \checkmark & \checkmark & \checkmark & \textbf{22.0} & 10.7 & 7.4 & 1.4 & \textbf{10.4} \\
\bottomrule
\end{tabular*}
\end{table*}

\subsubsection{Financial Reasoning: Formula and FiNER}

Table~\ref{tab:agentic-finance-ablation} reports an incremental ablation
that adds the three mechanisms on top of the Origin Playbook, which
obtains $70.0$ on Formula and $51.3$ on FiNER (average $60.7$). Adding RAE
raises Formula to $77.5$ by supplying the Generator with a compact set of
highly relevant rules; FiNER dips to $47.6$, indicating that it can
instead benefit from broader complementary context. Layering the
adversarial curriculum on top yields the strongest financial-reasoning
configuration, lifting Formula to $79.5$ and recovering FiNER to $52.1$
for the best average ($65.8$): boundary-case stress tests expose rule
gaps that routine training misses.

Adding the Failure Memory Bank completes the full deployed configuration.
It trades a little Formula accuracy ($73.0$) for the best FiNER score
($52.9$), reflecting FMB's role in supplying analogous, Ground
Truth-backed failures that stabilize cross-task behavior---an effect that
matters most on the interactive AppWorld agents examined next, where the
same failures recur across multi-step trajectories.

\subsubsection{Interactive Agents: AppWorld}

Table~\ref{tab:agentic-appworld-ablation} follows the same cumulative
order on AppWorld, using Task Goal Completion (TGC) and Scenario Goal
Completion (SGC) on the normal and challenge splits. Adding RAE lifts the
Origin average from $9.3$ to $10.4$, improving normal SGC ($8.9\to12.5$)
and challenge TGC ($6.7\to8.2$) by focusing the Generator on relevant
rules. The adversarial curriculum, valuable on the static financial
tasks, does not by itself transfer to interactive execution: layered on
RAE it improves only challenge SGC ($0.7\to2.2$) while the aggregate slips
to $8.9$, since a synthetically exposed weakness yields lasting benefit
only once it is diagnosed and retained.

Adding the Failure Memory Bank supplies exactly this retention: the full
deployed configuration recovers to the best average ($10.4$) and the best
normal TGC ($22.0$), confirming that verified failure recall is what makes
the adversarial signal durable when an agent must recover from recurrent
multi-step execution errors. Taken together, the two benchmark groups
tell a consistent story: RAE focuses execution, the adversarial curriculum
surfaces latent weaknesses, and the Failure Memory Bank retains their
lessons, so the full configuration is the most balanced across the static
and interactive tasks.

\subsection{End-to-End Efficiency: SCR on the Agentic Layer}
\label{subsec:exp-scr-agent}

The two preceding subsections evaluated each component on its own
bottleneck and baselines. We now measure them jointly, running the
self-improving agentic layer on top of the SCR inference engine---the
configuration DeepEdu actually deploys---to verify that the amortized
long-context engine translates into concrete latency gains for the
agent, whose every turn re-prefills a Playbook-conditioned context.

\paragraph{Setup}
We measure inference efficiency for ACE on the AppWorld normal and
challenge splits, using the same Playbook and Qwen3-4B-Instruct-2507
backbone on the H100 platform of §\ref{subsec:agentic-experiments}. The
\textbf{Origin} configuration is the standard ACE implementation served
through vLLM; \textbf{SCR} replaces its long-context execution path with
our Similarity Chunk Rolling engine (§\ref{subsec:long-context-engine})
while keeping the ACE task interface and splits unchanged. We report
time-to-first-token (TTFT), time-per-output-token (TPOT), and aggregate
input/output tokens.
\begin{figure}[H]
\centering
\includegraphics[width=0.96\columnwidth]{figures/ttft_origin.png}
\caption{Aggregate inference statistics for ACE on the AppWorld normal
and challenge splits}
\label{fig:ttft_origin}
\end{figure}

\paragraph{TTFT}
Figure~\ref{fig:ttft_origin} reports latency on the two AppWorld
splits. SCR reduces mean TTFT from $11.96$ to
$5.51$ seconds on Normal ($53.9\%$ reduction; $2.17\times$ faster)
and from $12.02$ to $5.75$ seconds on Challenge ($52.2\%$ reduction;
$2.09\times$ faster). This reduction is valuable for ACE because each
agent turn begins by prefilling a Playbook-conditioned context.



\paragraph{TPOT and token consumption}
SCR itself targets the prefill stage: the TTFT reductions above are its
direct effect, and it leaves the decode path unchanged. Per-token decode
latency (TPOT) is instead governed by the \emph{selection cache} of the
underlying TokenSelect engine~\cite{tokenselect2025}, which caches
token-selection results across consecutive, highly similar decode queries.
Because SCR is built on top of TokenSelect, the deployed engine retains
this cache and keeps TPOT low at decode; TokenSelect reports that the
mechanism lowers decode latency and yields up to a $2.28\times$ end-to-end
speedup. The two effects are complementary---SCR amortizes selection at
prefill, the selection cache amortizes it at decode---so the deployed
engine stays fast in both phases (Figure~\ref{fig:ttft_origin}).
% TODO: if reporting measured TPOT (Origin vs. deployed, seconds/token,
% Normal and Challenge) from the joint run, state the numbers here.
The deployed engine also processes fewer input tokens on both splits,
consistent with selective context execution; total output-token counts
differ across the evaluated runs, so they should not be interpreted as an
isolated throughput comparison.

\section{Conclusion}
\label{sec:conclusion}
We presented DeepEdu-v1, an AI-tutoring system for Vietnamese education
deployable on the constrained, on-premise hardware that data-sovereignty
regulations require. Its technical core is SCALE, a framework that couples
two components targeting the two bottlenecks such deployments face. \emph{Similarity Chunk Rolling} (SCR) attacks the
dynamic KV-cache and prefill-latency bottleneck by amortizing selective
sparse attention from per-sub-chunk to per-cluster granularity, sustaining
low latency as context grows. The \emph{self-improving agentic layer}
attacks the semantic bottleneck: rather than costly weight updates, it
continuously curates an evolving Playbook, using a failure-memory bank and
an adversarial curriculum to proactively expose and repair the systematic,
recurring weaknesses we observed in the base agent.

Our evaluation shows that these components are effective both separately
and together. On long-context benchmarks SCR preserves accuracy while
cutting prefill latency, and running the agentic layer on top of SCR
reduces mean time-to-first-token by roughly $2\times$ ($2.17\times$ on the
AppWorld normal split and $2.09\times$ on challenge)---the configuration
DeepEdu actually deploys---confirming that the two contributions compose
rather than merely coexist. Because SCR builds on TokenSelect, the
deployed engine also inherits its selection cache, which keeps per-token
decode latency low. The self-improving layer also lifts final accuracy on
complex tasks from $70.0\%$ to $79.5\%$ and yields the strongest
per-track gains across finer-grained sub-tasks and adversarial splits,
breaking the performance ceiling of the underlying framework without any
fine-tuning.

One main limitation frames our future work: our experiments validate the
enabling mechanisms on established long-context and agentic benchmarks,
whereas evaluating the localized grounding claim directly on native,
region-specific curricula---the setting that ultimately motivates
DeepEdu---remains important future work
toward democratizing trustworthy AI tutoring in resource-constrained
regions.


%\section*{Acknowledgments}
%This work was supported by the Posts and Telecommunications Institute of Technology.

\appendix
\section{Full InfiniteBench Configuration Grid}
\label{app:infbench-full}

Tables~\ref{tab:infbench-accuracy-full} and
\ref{tab:infbench-latency-full} report the complete results for all nine
SCR configurations ($\theta\in\{0.95,0.97,0.99\}$,
$L_{\max}\in\{1024,2048,4096\}$) across all six InfiniteBench tasks,
including R.Num and R.PK, which are
saturated at $100\%$ accuracy for every method. The representative
subsets in the main-text Tables~\ref{tab:infbench-accuracy} and
\ref{tab:infbench-latency} are drawn from these grids.

\begin{table*}[t]
\centering
\caption{Full per-task accuracy (\%) on InfiniteBench (Qwen2-7B-Instruct)
for all nine SCR configurations and the baselines. Best per column in
bold; \texttt{R.Num}/\texttt{R.PK} are saturated for the strong methods
and left unbolded.}
\label{tab:infbench-accuracy-full}
\small
\begin{tabular}{lcccccc}
\toprule
\textbf{Method} & \textbf{Code.D} & \textbf{R.KV} & \textbf{En.Dia} & \textbf{Math.F} & \textbf{R.Num} & \textbf{R.PK} \\
\midrule
Qwen2-7B~\cite{yang2024qwen2}          & 28.2 & 19.0 & \phantom{0}8.5 & 19.7 & 28.6 & 28.8 \\
NTK~\cite{ntk2023web}                  & 24.9 & 59.8 & \phantom{0}7.5 & \textbf{27.7} & 97.5 & 99.2 \\
SelfExtend~\cite{jin2024selfextend}    & \phantom{0}8.1 & \phantom{0}0.0 & \phantom{0}5.0 & \phantom{0}2.3 & \phantom{0}0.0 & \phantom{0}0.0 \\
StreamingLLM~\cite{xiao2024efficient}  & 27.9 & \phantom{0}2.4 & \phantom{0}3.5 & 19.4 & \phantom{0}5.1 & \phantom{0}5.1 \\
InfLLM~\cite{xiao2024infllm}           & 27.4 & \phantom{0}5.4 & \phantom{0}7.5 & 24.0 & 72.2 & 70.3 \\
TokenSelect~\cite{tokenselect2025}     & 26.7 & 91.0 & 10.0 & 24.0 & 100.0 & 100.0 \\
\midrule
$\theta\!=\!0.95$, $L_{\max}\!=\!1024$      & 27.6 & 94.0 & 10.0 & 23.7 & 100.0 & 100.0 \\
$\theta\!=\!0.95$, $L_{\max}\!=\!2048$      & 26.1 & 96.2 & 10.0 & 23.4 & 100.0 & 100.0 \\
$\theta\!=\!0.95$, $L_{\max}\!=\!4096$      & 24.9 & \textbf{98.0} & \phantom{0}6.0 & 22.6 & 100.0 & 100.0 \\
$\theta\!=\!0.97$, $L_{\max}\!=\!1024$      & 27.0 & 94.0 & 11.0 & 23.1 & 100.0 & 100.0 \\
$\theta\!=\!0.97$, $L_{\max}\!=\!2048$      & 27.9 & 96.2 & \phantom{0}9.0 & 22.9 & 100.0 & 100.0 \\
$\theta\!=\!0.97$, $L_{\max}\!=\!4096$      & 25.8 & \textbf{98.0} & \phantom{0}8.0 & 22.6 & 100.0 & 100.0 \\
$\theta\!=\!0.99$, $L_{\max}\!=\!1024$      & \textbf{28.8} & 94.0 & \textbf{11.5} & 23.1 & 100.0 & 100.0 \\
$\theta\!=\!0.99$, $L_{\max}\!=\!2048$      & 26.7 & 96.6 & 10.0 & 23.7 & 100.0 & 100.0 \\
$\theta\!=\!0.99$, $L_{\max}\!=\!4096$      & 25.8 & \textbf{98.0} & \phantom{0}7.0 & 23.4 & 100.0 & 100.0 \\
\bottomrule
\end{tabular}
\end{table*}

\begin{table*}[t]
\centering
\caption{Full per-task TTFT latency (s) on InfiniteBench for all nine SCR
configurations. Best (lowest) per column in bold.}
\label{tab:infbench-latency-full}
\small
\begin{tabular}{lcccccc}
\toprule
\textbf{Method} & \textbf{Code.D} & \textbf{R.KV} & \textbf{En.Dia} & \textbf{Math.F} & \textbf{R.Num} & \textbf{R.PK} \\
\midrule
TokenSelect                            & 8.8 & 13.2 & 7.7 & 8.7 & 9.4 & 9.3 \\
\midrule
$\theta\!=\!0.95$, $L_{\max}\!=\!1024$      & 6.7 & 9.9 & 6.0 & 6.6 & 7.1 & 7.0 \\
$\theta\!=\!0.95$, $L_{\max}\!=\!2048$      & 5.9 & 8.7 & 5.3 & 5.9 & 6.3 & 6.3 \\
$\theta\!=\!0.95$, $L_{\max}\!=\!4096$      & \textbf{5.7} & \textbf{8.3} & \textbf{5.1} & \textbf{5.7} & \textbf{6.1} & \textbf{6.0} \\
$\theta\!=\!0.97$, $L_{\max}\!=\!1024$      & 6.7 & 9.8 & 6.0 & 6.6 & 7.0 & 7.0 \\
$\theta\!=\!0.97$, $L_{\max}\!=\!2048$      & 6.1 & 8.7 & 5.4 & 5.9 & 6.4 & 6.4 \\
$\theta\!=\!0.97$, $L_{\max}\!=\!4096$      & 5.8 & \textbf{8.3} & \textbf{5.1} & \textbf{5.7} & \textbf{6.1} & 6.1 \\
$\theta\!=\!0.99$, $L_{\max}\!=\!1024$      & 7.3 & 9.9 & 6.2 & 6.7 & 7.4 & 7.5 \\
$\theta\!=\!0.99$, $L_{\max}\!=\!2048$      & 6.8 & 8.7 & 5.7 & 5.9 & 6.7 & 6.6 \\
$\theta\!=\!0.99$, $L_{\max}\!=\!4096$      & 6.7 & \textbf{8.3} & 5.7 & \textbf{5.7} & 6.3 & 6.3 \\
\bottomrule
\end{tabular}
\end{table*}

\bibliography{references}

\end{document}